% Exacta user manual. Build with docs/manual/build.sh (or: latexmk -pdf userManual.tex).
% Chapters adapted from the Wiki pages in docs/wiki are marked at the top of each file.
\documentclass[11pt,a4paper]{report}

\usepackage[T1]{fontenc}
\usepackage[utf8]{inputenc}
\usepackage{lmodern}
\usepackage{textcomp}
\usepackage[margin=2.6cm]{geometry}
\usepackage{microtype}
\usepackage{amsmath,amssymb}
\usepackage{calc}
\usepackage{array,booktabs,longtable}
\usepackage{etoolbox}
\usepackage{enumitem}
\usepackage{xcolor}
\usepackage{listings}
\usepackage[hidelinks,pdftitle={Exacta User Manual},pdfauthor={Rodolfo Orlando Viana}]{hyperref}
\usepackage{xurl}

% The version printed on the title page. build.sh writes version.tex from src/exacta/__init__.py;
% without it (a plain latexmk run) the manual says which release it was last built for.
\IfFileExists{version.tex}{\input{version.tex}}{\newcommand{\exactaversion}{0.1.3}}

\emergencystretch=3em
% Let long identifiers in typewriter type break at the end of a line (TeX hyphenates them).
\DeclareRobustCommand{\texttt}[1]{{\normalfont\ttfamily\hyphenchar\font=`\-\relax #1}}
\let\textunderscoreorig\_
\renewcommand{\_}{\textunderscoreorig\allowbreak}
\setlength{\parindent}{0pt}
\setlength{\parskip}{0.55\baselineskip plus 2pt}
\setcounter{secnumdepth}{2}
\setcounter{tocdepth}{2}
\providecommand{\tightlist}{\setlength{\itemsep}{0pt}\setlength{\parskip}{0pt}}
\newcounter{none} % for the unnumbered tables that pandoc writes (\LTcaptype{none})
\makeatletter
\patchcmd\longtable{\par}{\if@noskipsec\mbox{}\fi\par}{}{}
\makeatother

\definecolor{codebg}{gray}{0.96}
\lstset{
  basicstyle=\ttfamily\footnotesize,
  backgroundcolor=\color{codebg},
  breaklines=true,
  breakatwhitespace=false,
  columns=fullflexible,
  keepspaces=true,
  upquote=true,
  showstringspaces=false,
  frame=single,
  framerule=0pt,
  xleftmargin=4pt,
  xrightmargin=4pt,
  aboveskip=0.6\baselineskip,
  belowskip=0.3\baselineskip,
  keywordstyle=\bfseries,
  commentstyle=\itshape\color{black!65},
}
\lstdefinelanguage{output}{}  % program output: no highlighting

\newcommand{\exacta}{\textsc{Exacta}}

\title{\exacta{} User Manual\\[0.4em]
  \large Fast, reference-validated losses and metrics\\ for deep learning on CPU and GPU\\[0.8em]
  Version \exactaversion}
\author{Rodolfo Orlando Viana\\[0.3em]
  \normalsize Institute of Biosciences, Humanities and Exact Sciences (IBILCE),\\
  \normalsize S\~ao Paulo State University (UNESP), S\~ao Jos\'e do Rio Preto, SP 15054-000, Brazil\\
  \normalsize\texttt{rodolfo.viana@unesp.br}}
\date{}

\begin{document}
\maketitle

\begin{abstract}
\noindent
Exacta is a Python library, with C++ and Triton kernels, for evaluating and training image
segmentation models in two and three dimensions. This manual explains how to install it, how
to compute its metrics and losses, which conventions each score follows, how to reproduce the
definitions of other tools, and how to run the example driver and the referee test subset that
come with the source. It is written for users and for readers who want to check the code.
\end{abstract}

\paragraph{Acknowledgements.}
This study was financed in part by the Coordena\c{c}\~ao de Aperfei\c{c}oamento de Pessoal de
N\'ivel Superior -- Brasil (CAPES) -- Finance Code 001.

\tableofcontents

\chapter{Introduction}\label{ch:introduction}

Exacta provides metrics and differentiable losses for two- and three-dimensional image
segmentation, with a focus on medical images. Overlap scores (Dice, IoU, sensitivity and
others), surface-distance scores (Hausdorff distance, HD95, average symmetric surface distance
and normalized surface Dice), physical areas and volumes, probability and calibration scores,
lesion and instance evaluation, dataset-level classification and regression metrics, and the
usual training objectives (Dice, cross-entropy, focal, Tversky, boundary and topology losses)
share one interface. CPU work runs in C++ kernels with OpenMP; supported CUDA work runs in
Triton kernels or in native PyTorch.

\section{Why conventions are arguments}\label{introduction:why}

Two tools can report different numbers for a score that has the same name. HD95 may be the
maximum of the two directional 95th percentiles or one percentile over both directions
together; ASSD may weight every surface voxel equally or by surface area; NSD may count a
distance that equals the tolerance, or not, depending on how the last bit of the distance
rounds. Exacta makes each of these choices an argument with a documented default, so that a
score is defined by its name and its arguments. Chapter~\ref{ch:conventions} states every
convention, Chapter~\ref{ch:tools} gives the arguments that reproduce MONAI, MedPy, DeepMind's
surface-distance and MeshMetrics, and the \emph{reporting signature}
(Section~\ref{metrics:reporting-signatures-since-0.1.0}) writes the full definition of a score in one string to
print next to the number.

\section{What is guaranteed}\label{introduction:guarantees}

The accuracy claims differ between components. The table separates what is argued or checked
exactly from what holds to floating-point rounding. The last column names the test files that
check each claim; Chapter~\ref{ch:tests} says how to run them.

\begin{longtable}{@{}
  >{\raggedright\arraybackslash}p{0.22\linewidth}
  >{\raggedright\arraybackslash}p{0.48\linewidth}
  >{\raggedright\arraybackslash}p{0.24\linewidth}@{}}
\toprule
Component & Guarantee & Checked by \\
\midrule
\endhead
Squared Euclidean distance transform &
The result is the exact floating-point minimum, over the feature voxels, of the sum over the axes,
in axis order, of $\mathrm{fl}\big(\mathrm{fl}((q_i-p_i)h_i)^2\big)$ with every product and sum
rounded once. The C++, Triton and PyTorch implementations return the same bits. &
\texttt{test\_edt.py}: brute-force oracle, property test, lines found from the exactness argument \\
Hausdorff distance, HD95 &
Exact for the discrete definition (the distances between boundary voxels, surface elements or
triangle centres), not for the continuous shape that the mask samples. &
\texttt{test\_surface.py}, \texttt{test\_surfels.py}, \texttt{test\_mesh.py} \\
ASSD, weighted NSD &
Sums over surface elements are correctly rounded (the exact sum rounded once), so the bits do
not depend on element order, thread count or device. &
\texttt{test\_sums.py} \\
Mesh distances &
Distances from triangle centres to the other mesh are minima of binary64 point--triangle
distances over all elements. This is tested against brute force, not proved. &
\texttt{test\_mesh.py} \\
Overlap and volume scores &
Computed from exact integer counts of true positives, false positives and false negatives. &
\texttt{test\_overlap.py}, \texttt{test\_volume.py} \\
Losses &
Precise to the rounding of the stated formula in the stated precision. Their bits can differ
between PyTorch builds and between CPU and GPU. &
\texttt{test\_losses.py}, \texttt{test\_composed\_losses.py} and others \\
\bottomrule
\end{longtable}

An exact score is a property of the discrete definition. Exacta does not claim that a metric
computed on a voxel grid equals the metric of the continuous shapes that the masks represent,
and it does not claim that any loss improves a model. A convergence of tools on one number
needs the same definition and an exact computation of it; Exacta supplies the second and makes
the first explicit.

\section{How to read this manual}\label{introduction:reading}

Chapters~\ref{ch:installation} and~\ref{ch:quickstart} install the library and run a first
example. Chapter~\ref{ch:conventions} defines the surface scores and the policies for empty
masks and ties. Chapters~\ref{ch:metrics} and~\ref{ch:losses} describe the metrics and the
losses; Chapter~\ref{ch:tools} lists the arguments that reproduce other tools, and
Chapter~\ref{ch:performance} reports measured speed. Chapter~\ref{ch:other} lists the modules
that this manual does not describe in detail. Chapters~\ref{ch:driver} and~\ref{ch:tests}
describe the example driver and the test subset for referees, and Chapter~\ref{ch:citing}
says how to cite the software.

Code, argument names and file names are set in \texttt{typewriter}. Spatial arrays have shape
$(B,\ast\text{spatial})$ for class indices and $(B,C,\ast\text{spatial})$ for channels, with $B$
cases and $C$ classes. $P$ is the prediction and $T$ the target. Spacings are in the axis
order of the array, in the unit of the reported distances (millimetres in the examples).

% Adapted from docs/wiki/Installation.md; edit this file, not the wiki page, for the manual.
\chapter{Installation}\label{ch:installation}

Exacta requires Python 3.10 or newer, PyTorch 2.4 or newer and NumPy 1.22 or newer.
Install PyTorch for your hardware using its \href{https://pytorch.org/get-started/locally/}{installation guide},
then install Exacta. Pin the release you evaluated, so that results can be traced to it:

\begin{lstlisting}[language=bash]
python -m pip install exacta==X.Y.Z   # X.Y.Z: the release you use
python -c "import exacta; print(exacta.__version__)"
\end{lstlisting}

The wheels are for Linux x86-64 with glibc 2.28 or newer. Python 3.10 and 3.11
have dedicated wheels; the Python 3.12 stable-ABI wheel supports Python 3.12 and newer.

Other platforms install from source, which builds the CPU extension and requires a C++17
compiler and Python development headers; OpenMP is optional. On Debian/Ubuntu, install
\texttt{build-essential} and development headers matching your Python interpreter.

\section{CPU and GPU}\label{installation:cpu-and-gpu}

CPU kernels are precompiled in wheels and compile during source installation. CUDA loss and
metric kernels use Triton at runtime; their first invocation can include compilation. The CPU extension does not link against a
specific PyTorch binary. Local validation covers Linux x86-64 CPU and NVIDIA CUDA execution;
other hardware/platform combinations should be verified before use.

\section{Distributed CPU execution}\label{installation:distributed-cpu-execution}

For Gloo distributed runs, use PyTorch 2.8 or newer. PyTorch 2.4.1 can deadlock during
process-group destruction after collecting metric scores: a Gloo worker needs the Python
GIL while the outer destructor holds it. PyTorch's \href{https://github.com/pytorch/pytorch/releases/tag/v2.8.0}{2.8 release notes}
include the destructor fix (\#154976). Ordinary non-distributed CPU evaluation still supports
PyTorch 2.4.

\section{Optional references}\label{installation:optional-references}

SciPy is needed for Hungarian lesion matching. MONAI, MedPy, VTK, SimpleITK and other reference
libraries are used by optional validation tests; they are not needed for ordinary Dice or
HD95 evaluation. Chapter~\ref{ch:tests} lists the packages and how to install them.

\section{Development checkout}\label{installation:development-checkout}

Use \texttt{main} to evaluate merged changes beyond the published release. Create a feature branch
from it when contributing. This manual describes the version printed on its title page; a newer
checkout may expose additional APIs or changed behavior. To keep your own Python environment
untouched, create a virtual environment first (Section~\ref{tests:referee}).

\begin{lstlisting}[language=bash]
git clone --branch main https://github.com/rodolfo-viana/exacta.git
cd exacta
python -m pip install -e ".[test]"
\end{lstlisting}

Record the commit SHA and environment alongside experiment results. Reinstall after changing
C++ sources so the compiled extension matches the checkout.

% Adapted from docs/wiki/Quickstart.md; edit this file, not the wiki page, for the manual.
\chapter{Quickstart}\label{ch:quickstart}

This synthetic example runs one training loss and evaluates hard labels. It fixes the random
seed and uses small 3D arrays; it is an API demonstration, not a segmentation experiment.

\begin{lstlisting}[language=Python]
import torch
import exacta.functional as EF
from exacta.losses import DiceCELoss
from exacta.metrics import Dice, HD95

generator = torch.Generator().manual_seed(0)
device = torch.device("cuda" if torch.cuda.is_available() else "cpu")
logits = torch.randn(2, 3, 8, 8, 8, generator=generator).to(device).requires_grad_()
target = torch.randint(3, (2, 8, 8, 8), generator=generator).to(device)

loss = DiceCELoss(include_background=False)(logits, target)
loss.backward()
assert torch.isfinite(loss) and torch.isfinite(logits.grad).all()

prediction = logits.detach().argmax(1)
spacing = (1.5, 0.8, 0.8)  # axis order of the array, in millimetres
per_case = EF.hd95(prediction, target, spacing=spacing, num_classes=3,
                   include_background=False)
assert per_case.shape == (2, 2)  # two cases and two foreground classes

dice = Dice(num_classes=3, include_background=False)
hd95 = HD95(num_classes=3, include_background=False)
dice.update(prediction, target)
hd95.update(prediction, target, spacing=spacing)
result = hd95.compute()
print("Loss:", loss.item())
print("HD95 per case:", result.per_case)
print("HD95 per class:", result.per_class)
print("HD95 mean:", result.mean)
\end{lstlisting}

Use \texttt{reset()} on a stateful metric at the start of a new evaluation. Metrics consume hard
labels; losses consume logits unless an overlap loss explicitly uses \texttt{activation="none"}.
Review Conventions (Chapter~\ref{ch:conventions}) before comparing
results against another package.

% Adapted from docs/wiki/Conventions.md; edit this file, not the wiki page, for the manual.
\chapter{Mathematical conventions}\label{ch:conventions}

Reference libraries sometimes implement different quantities under the same name. Exacta
makes these choices arguments; matching a library requires matching all relevant settings.

{\small\def\LTcaptype{none} % do not increment counter
\begin{longtable}[]{@{}
  >{\raggedright\arraybackslash}p{(\linewidth - 2\tabcolsep) * \real{0.2700}}
  >{\raggedright\arraybackslash}p{(\linewidth - 2\tabcolsep) * \real{0.7300}}@{}}
\toprule\noalign{}
\begin{minipage}[b]{\linewidth}\raggedright
Choice
\end{minipage} & \begin{minipage}[b]{\linewidth}\raggedright
Exacta behavior
\end{minipage} \\
\midrule\noalign{}
\endhead
\bottomrule\noalign{}
\endlastfoot
HD95 directions & \texttt{combine="max"}: maximum of directional percentiles; \texttt{"pooled"}: percentile over both directions together \\
ASSD directions & \texttt{combine="pooled"}: mean over both sets of distances; \texttt{"mean"}: mean of the two directional means \\
Voxel surfaces & \texttt{surface="voxels"}: foreground boundary voxels with equal weights; face-adjacent background identifies the boundary \\
Surfels & \texttt{surface="surfels"}: marching-cubes surface elements weighted by area (contour lengths in 2D); percentiles use these weights \\
Meshes & \texttt{surface="mesh"}: triangle centers to nearest points on the other mesh, weighted by triangle area (contour segments in 2D); evaluated on CPU \\
NSD threshold & Distances at or below tolerance are included, up to a relative margin \texttt{tie\_rtol} (default \texttt{1e-6}, since 0.1.0) that keeps exact ties within despite rounding; \texttt{tie\_rtol=0} compares strictly \\
Sums over surface elements & Means (ASSD) and weighted NSD fractions use correctly rounded sums (since 0.1.2), so their bits do not depend on element order, thread count or device \\
Image exterior & Counts as background \\
Both masks empty & \texttt{both\_empty}, default NaN \\
One mask empty & Dice/IoU: zero; HD/ASSD: \texttt{one\_empty}, default infinity; NSD: \texttt{one\_empty}, default zero \\
Undefined ratio & \texttt{zero\_division}, default NaN, for precision/sensitivity/specificity and volume ratios \\
\end{longtable}
}

\section{Surface equations}\label{conventions:surface-equations}

For two nonempty surfaces, let \(d_i^P\) denote the distance from prediction surface element
\(i\) to the target surface, and \(d_j^T\) the reverse distance. Distances use physical coordinates
from \texttt{spacing}. Weights \(w_i^P,w_j^T\) are one for voxel surfaces and area (length in 2D) for
surfels or meshes, using the distance geometry described above.

For \texttt{combine="max"}, HD95 takes the larger directional 95th percentile:

\[
\mathrm{HD95}_{\mathrm{max}}
=\max\left(Q_{0.95}(d^P;w^P),Q_{0.95}(d^T;w^T)\right).
\]

\(Q_{0.95}\) uses linear interpolation for voxel distances and an area-weighted percentile
for surfels and meshes. \texttt{combine="pooled"} takes one percentile after concatenating the two
distance/weight sequences. These two operations need not agree.

The default pooled ASSD weights all surface elements together:

\[
\mathrm{ASSD}_{\mathrm{pooled}}
=\frac{\sum_i w_i^P d_i^P+\sum_j w_j^T d_j^T}
{\sum_i w_i^P+\sum_j w_j^T}.
\]

For \texttt{combine="mean"}, take each direction's weighted mean separately, then average the two
means with equal weight. NSD at tolerance \(\tau\) is the weighted fraction within tolerance,
where \(\mathbf{1}[\cdot]\) is one when the condition holds and zero otherwise:

\[
\mathrm{NSD}_{\tau}
=\frac{\sum_i w_i^P\mathbf{1}[d_i^P\leq\tau]
+\sum_j w_j^T\mathbf{1}[d_j^T\leq\tau]}
{\sum_i w_i^P+\sum_j w_j^T}.
\]

\textbf{Ties with the tolerance (since 0.1.0).} A distance and a tolerance that are equal in exact
arithmetic can differ after rounding. With spacing 0.4 mm, three voxels evaluate to
\(\mathrm{fl}(3\times0.4)=1.2000000000000002\) in binary64, above a tolerance of 1.2 mm, and a
spacing read from a float32 image header carries a relative error up to \(6\times10^{-8}\). A strict
comparison then excludes every surface element lying exactly on the tolerance, which can move NSD
by several percentage points. Exacta therefore evaluates \(\mathbf{1}[d\leq\tau(1+\rho)]\), where
\(\rho\) is \texttt{tie\_rtol} (default \(10^{-6}\), at most \(10^{-3}\)). The default exceeds both rounding
errors and stays below the relative gap between distinct distances on an isotropic grid, about
\(1/(2n)\) for a squared distance of \(n\) voxels, for tolerances up to about 700 voxels. It never
widens a tolerance of zero. \texttt{tie\_rtol=0} reproduces SciPy-based tools (MedPy, DeepMind's
surface-distance), which compare in binary64; MONAI compares float32 distances, which counts
exact ties as the default does.

Empty surfaces use the explicit policies in the table rather than these ratios.

\section{Reproducing reference definitions}\label{conventions:reproducing-reference-definitions}

MedPy's HD95 pools directions; MONAI's usual voxel HD95 takes their maximum. DeepMind's
surface-distance uses surfel areas. MeshMetrics uses mesh geometry. ASSD changed between
MedPy 0.4 (mean of directional means) and 0.5.2 (pooled mean). Record library versions as well
as Exacta's arguments when reporting comparisons.

With \texttt{volume\_similarity(convention="taha")} the score is \texttt{1\ -\ abs(Vp-Vt)/(Vp+Vt)}.
With \texttt{convention="itk"} it is the signed \texttt{2*(Vp-Vt)/(Vp+Vt)}. These are different quantities.

\section{Physical axes and numerical limits}\label{conventions:physical-axes-and-numerical-limits}

Spacing follows the array's axes. A SimpleITK array commonly uses \texttt{(z,\ y,\ x)}, whereas
\texttt{GetSpacing()} returns \texttt{(x,\ y,\ z)}; reverse it when those are the layouts in use.
All distances and NSD tolerances share the spacing unit.

Voxel-surface spacing components must be in \texttt{{[}1e-100,\ 1e100{]}}; surfels and meshes require
\texttt{{[}1e-50,\ 1e50{]}} because their geometric calculations include higher powers of coordinates.
Unsupported spacing raises even for empty masks. Small floating-point differences near
an NSD threshold can change membership; do not silently replace reference precision or
enable compiler transformations that change rounding.

% Adapted from docs/wiki/Metrics.md; edit this file, not the wiki page, for the manual.
\chapter{Metrics}\label{ch:metrics}

Overlap and surface functions in \texttt{exacta.functional} return per-case, per-selected-class
arrays of shape \texttt{(B,\ K)}. Corresponding classes in \texttt{exacta.metrics} accumulate results across
batches. Probability and dataset-level metrics use the reductions described below.
Choose complementary metrics according to the segmentation task; a high overlap score alone
does not describe lesion detection or boundary errors.

{\small\def\LTcaptype{none} % do not increment counter
\begin{longtable}[]{@{}
  >{\raggedright\arraybackslash}p{(\linewidth - 4\tabcolsep) * \real{0.2400}}
  >{\raggedright\arraybackslash}p{(\linewidth - 4\tabcolsep) * \real{0.3000}}
  >{\raggedright\arraybackslash}p{(\linewidth - 4\tabcolsep) * \real{0.4600}}@{}}
\toprule\noalign{}
\begin{minipage}[b]{\linewidth}\raggedright
Functional name
\end{minipage} & \begin{minipage}[b]{\linewidth}\raggedright
Stateful class
\end{minipage} & \begin{minipage}[b]{\linewidth}\raggedright
Definition
\end{minipage} \\
\midrule\noalign{}
\endhead
\bottomrule\noalign{}
\endlastfoot
\texttt{dice} & \texttt{Dice} & \(2\mathrm{TP}/(2\mathrm{TP}+\mathrm{FP}+\mathrm{FN})\) \\
\texttt{iou} & \texttt{IoU} & \(\mathrm{TP}/(\mathrm{TP}+\mathrm{FP}+\mathrm{FN})\) \\
\texttt{fbeta} & \texttt{FBeta} & \((1+\beta^2)\mathrm{TP}/\big((1+\beta^2)\mathrm{TP}+\beta^2\mathrm{FN}+\mathrm{FP}\big)\); \(\beta=1\) is Dice \\
\texttt{sensitivity} & \texttt{Sensitivity} & \(\mathrm{TP}/(\mathrm{TP}+\mathrm{FN})\) \\
\texttt{precision} & \texttt{Precision} & \(\mathrm{TP}/(\mathrm{TP}+\mathrm{FP})\) \\
\texttt{specificity} & \texttt{Specificity} & \(\mathrm{TN}/(\mathrm{TN}+\mathrm{FP})\) \\
\texttt{volume\_similarity} & \texttt{VolumeSimilarity} & Default: \(1-\lvert V_P-V_T\rvert/(V_P+V_T)\) \\
\texttt{relative\_volume\_difference} & \texttt{RelativeVolumeDifference} & \((V_P-V_T)/V_T\); \texttt{absolute=True} removes the sign \\
\texttt{hausdorff\_distance} & \texttt{HausdorffDistance} & Percentile boundary distance; default percentile 100 \\
\texttt{hd95} & \texttt{HD95} & 95th-percentile boundary distance \\
\texttt{assd} & \texttt{ASSD} & Average symmetric boundary distance \\
\texttt{nsd} & \texttt{NSD} & Fraction of both surfaces within the supplied tolerance \\
\end{longtable}
}

For each case and class, TP, FP, FN and TN count true-positive, false-positive,
false-negative and true-negative voxels. \(V_P\) and \(V_T\) are predicted and target foreground
voxel counts. The formulas apply when their denominators are nonzero; empty-mask and
zero-denominator policies are specified in
Conventions (Chapter~\ref{ch:conventions}).

Inputs use class indices \texttt{(B,\ *spatial)} with \texttt{num\_classes}, or binary masks \texttt{(B,\ C,\ *spatial)}
with \texttt{input\_format="one-hot"}. \texttt{include\_background=False} omits channel/class zero.
Stateless functions accept PyTorch, NumPy, CuPy and JAX arrays and return the prediction's
array type. Stateful metrics store PyTorch results.

\section{Selected voxels and ignored labels}\label{metrics:selected-voxels-and-ignored-labels}

The seven overlap/volume metrics listed above accept \texttt{mask} and \texttt{ignore\_index}.
For index inputs, \texttt{ignore\_index} excludes voxels whose \textbf{target} equals that label. It does
not remove a class from the output. A boolean \texttt{mask} includes True voxels and excludes False
voxels; its shape is \texttt{(B,\ *spatial)} or \texttt{(B,\ 1,\ *spatial)} on the input device. For one-hot or
multilabel masks, the same voxel selection applies to every channel; use \texttt{mask} instead of
\texttt{ignore\_index}. When both options are supplied, a voxel must pass both conditions.

Let \(V\) be the included voxel set. All four counts are restricted to \(V\):

\[
\mathrm{TP}=\lvert P\cap T\cap V\rvert,\quad
\mathrm{FP}=\lvert(P\setminus T)\cap V\rvert,\quad
\mathrm{FN}=\lvert(T\setminus P)\cap V\rvert,\quad
\mathrm{TN}=\lvert V\setminus(P\cup T)\rvert.
\]

Excluded values in either input are not validated or counted. An entirely excluded case has
zero counts and follows \texttt{both\_empty} or \texttt{zero\_division}, depending on the metric. The default
NaN is skipped by case averaging; choosing a finite fill explicitly includes it in that average.
Masks are consumed per call and are not retained by stateful metrics.

For example, evaluate synthetic partially annotated labels:

\begin{lstlisting}[language=Python]
import torch
from exacta import functional as EF
from exacta.metrics import MetricCollection

target = torch.tensor([[[1, 1], [0, -1]]])
prediction = torch.tensor([[[1, 0], [1, 999]]])
roi = torch.tensor([[[True, True], [True, True]]])
scores = EF.dice(prediction, target, num_classes=2, ignore_index=-1, mask=roi)
assert torch.allclose(scores, torch.tensor([[0., .5]], dtype=torch.float64))
metrics = MetricCollection(metrics=["dice", "specificity"], num_classes=2, ignore_index=-1)
metrics.update(prediction, target, mask=roi)
results = metrics.compute()
\end{lstlisting}

The CPU C++ and CUDA Triton paths exclude voxels inside the count pass, without one-hot
expansion of index maps. A checked PyTorch fallback covers other devices and larger class
histograms. Collections with surface or clDice metrics reject \texttt{mask} and \texttt{ignore\_index}:
turning excluded regions into background would create artificial boundaries or alter topology.

\section{Spacing and tolerances}\label{metrics:spacing-and-tolerances}

Distance metrics require spacing in the tensor's spatial-axis order. Supply one spatial
tuple for the batch or a \texttt{(B,\ ndim)} array for per-case spacing. The distance unit is the
spacing unit. NSD requires \texttt{tolerance}, which may be scalar or per selected class. Distances
equal to the tolerance count as within it even after rounding, through the relative margin
\texttt{tie\_rtol} (default \texttt{1e-6}, since 0.1.0); see
Conventions (Chapter~\ref{ch:conventions}).

For an example using synthetic hard labels:

\begin{lstlisting}[language=Python]
import numpy as np
import exacta.functional as EF

target = np.zeros((1, 8, 8, 8), dtype=np.int64)
target[:, 2:5, 2:5, 2:5] = 1
prediction = np.roll(target, 1, axis=1)
scores = EF.nsd(prediction, target, num_classes=2, include_background=False,
                spacing=(1.5, 0.8, 0.8), tolerance=1.0, surface="surfels")
assert scores.shape == (1, 1)
\end{lstlisting}

\section{Reporting signatures (since 0.1.0)}\label{metrics:reporting-signatures-since-0.1.0}

A metric's name does not fix its value: HD95 depends on the surface representation and on how
the two directions are joined, and NSD on how distances equal to the tolerance are compared.
\texttt{exacta.conventions.signature} writes every such choice, with defaults filled in, and the Exacta
version, in one string to report next to the number:

\begin{lstlisting}[language=Python]
from exacta.conventions import parse_signature, signature

text = signature("nsd", tolerance=2.0)
# 'nsd[both_empty=nan;one_empty=0.0;surface=voxels;tie_rtol=1e-06;tolerance=2.0]@exacta-0.1.2'
name, options, version = parse_signature(text)  # options for functional.evaluate
\end{lstlisting}

\texttt{parse\_signature} returns options that make \texttt{functional.evaluate} recompute the same score bit
for bit. Floats are written as Python's \texttt{repr}, so \texttt{nan}, \texttt{inf} and values such as \texttt{1e-06} read
back exactly. \texttt{signatures(selection)} covers a selection, and \texttt{MetricCollection.signatures()}
adds \texttt{aggregate=cases} or \texttt{aggregate=pooled}, which \texttt{parse\_signature} returns in the options and
\texttt{evaluate} does not accept; pop it before calling \texttt{evaluate}.

\section{Aggregation and distributed evaluation}\label{metrics:aggregation-and-distributed-evaluation}

\texttt{compute()} returns \texttt{per\_case}, \texttt{per\_class} and \texttt{mean}. Aggregation averages cases with NaNs
excluded, then averages the class means. It does not pool all voxels across patients.
In distributed evaluation every rank must call \texttt{compute()}, including ranks with no local
updates. With NCCL, set each rank's current CUDA device. Avoid samplers that pad by duplicating
cases. If no rank has any cases, \texttt{compute()} raises \texttt{RuntimeError}.

\subsection{Dataset-pooled counts}\label{metrics:dataset-pooled-counts}

\texttt{Dice}, \texttt{IoU}, \texttt{FBeta}, \texttt{Sensitivity}, \texttt{Precision}, \texttt{Specificity}, \texttt{VolumeSimilarity},
\texttt{RelativeVolumeDifference} and \texttt{MetricCollection} accept \texttt{aggregate="pooled"}. They then keep
only the totals \(\mathrm{TP}_c=\sum_n \mathrm{TP}_{n,c}\) (and FP, FN, TN) over every case \(n\)
and process, an \(O(C)\) state, and score those totals once per class \(c\). With
\(\mathcal{C}\) the selected classes whose pooled score is defined, the mean is

\[
\mathrm{mIoU}=\frac{1}{\lvert\mathcal{C}\rvert}\sum_{c\in\mathcal{C}}
\frac{\mathrm{TP}_c}{\mathrm{TP}_c+\mathrm{FP}_c+\mathrm{FN}_c},
\]

mmsegmentation's \texttt{IoUMetric} convention: a class absent from every prediction and target is
NaN and left out, while a class that is only ever predicted scores 0 and counts. Pooled Dice
is the ``global Dice'' some challenges report; the class mean of pooled sensitivity is
mmsegmentation's mean accuracy. \texttt{per\_case} is \texttt{None} in this mode.

\begin{lstlisting}[language=Python]
from exacta.metrics import IoU

miou = IoU(num_classes=19, ignore_index=255, aggregate="pooled")
for prediction, target in loader:
    miou.update(prediction, target)
result = miou.compute()  # result.per_class: IoU of each class; result.mean: mIoU
\end{lstlisting}

Pooling weights every voxel equally, so large structures and large images dominate; a case
mean weights every case equally. They are different estimands, and neither is a patient-level
3D score for 2D slices. Pooled collections accept only ratios of counts, not distances, clDice
or physical sizes. \texttt{functional.confusion\_counts} returns the \texttt{(B,\ K)} TP, FP, FN and TN arrays
behind every overlap metric, for definitions Exacta does not name or custom pooling such as one
patient's slices. Voxel accuracy and MCC depend on TN, and so on the field of view; dataset
pixel accuracy is \texttt{ClassificationMetrics(task="multiclass",\ metrics=("accuracy",))} applied to
flattened labels.

See Conventions (Chapter~\ref{ch:conventions}) for empty masks,
surface representations and reference-library differences. Excluding a region from a surface
metric needs its own definition; overlap masking does not define partial surface evaluation.

\section{Shared metric evaluation}\label{metrics:shared-metric-evaluation}

\texttt{functional.evaluate} shares confusion counts across overlap/area/volume metrics and
surface distances across compatible boundary metrics. It returns a dictionary of \texttt{(B,\ K)}
arrays. Choose names as a sequence, or map names to their individual options:

\begin{lstlisting}[language=Python]
selection = {"dice": {}, "hd95": {"combine": "pooled"}, "assd": {},
             "nsd": {"tolerance": 1.0, "surface": "voxels"}}
scores = EF.evaluate(prediction, target, metrics=selection, num_classes=2,
                     include_background=False, spacing=(1.5, 0.8, 0.8))
\end{lstlisting}

Each metric keeps its standalone defaults and empty-mask rules. Input layout, class selection
and spacing are shared; different surface representations form separate calculation groups.
Distances are reduced immediately per chunk, and no volume cache survives the call.

\texttt{MetricCollection} provides the same sharing during stateful evaluation, retaining only
detached scores and gathering them in one distributed payload:

\begin{lstlisting}[language=Python]
from exacta.metrics import MetricCollection

metrics = MetricCollection(metrics=selection, num_classes=2, include_background=False)
metrics.update(prediction, target, spacing=(1.5, 0.8, 0.8))
results = metrics.compute()
mean_hd95 = results["hd95"].mean
metrics.reset()
\end{lstlisting}

All ranks must use identical metric order and configuration, and must call \texttt{compute()},
including ranks without updates. The usual sampler and aggregation rules above still apply.

\section{Centreline, probability and detection metrics}\label{metrics:centreline-probability-and-detection-metrics}

The following APIs evaluate centreline overlap, probabilities and lesion detection. The
\href{https://github.com/rodolfo-viana/exacta/wiki/Advanced-segmentation}{advanced segmentation guide}
defines their empty-input behavior, probability conventions and matching protocol.

{\small\def\LTcaptype{none} % do not increment counter
\begin{longtable}[]{@{}
  >{\raggedright\arraybackslash}p{(\linewidth - 4\tabcolsep) * \real{0.3400}}
  >{\raggedright\arraybackslash}p{(\linewidth - 4\tabcolsep) * \real{0.3300}}
  >{\raggedright\arraybackslash}p{(\linewidth - 4\tabcolsep) * \real{0.3300}}@{}}
\toprule\noalign{}
\begin{minipage}[b]{\linewidth}\raggedright
Functional API
\end{minipage} & \begin{minipage}[b]{\linewidth}\raggedright
Stateful API
\end{minipage} & \begin{minipage}[b]{\linewidth}\raggedright
Quantity
\end{minipage} \\
\midrule\noalign{}
\endhead
\bottomrule\noalign{}
\endlastfoot
\texttt{cldice} & \texttt{ClDice} & Harmonic mean of hard-skeleton precision and sensitivity; CPU Zhang (2D) or Lee (3D) thinning \\
\texttt{betti\_error} & \texttt{BettiError} & \(\lvert\beta_k(P)-\beta_k(T)\rvert\): components, holes/tunnels or cavities \\
\texttt{brier\_score} & \texttt{BrierScore} & Mean squared probability error; binary and multiclass scales differ \\
\texttt{negative\_log\_likelihood} & \texttt{NegativeLogLikelihood} & Mean true-class negative log probability, in nats \\
\texttt{expected\_calibration\_error} & \texttt{ExpectedCalibrationError} & Top-label or classwise, equal-width reliability bins \\
\texttt{reliability\_diagram} & \texttt{Calibration} & Per-case summaries versus pooled voxel bin accumulation \\
\texttt{detection.lesion\_detection} & \texttt{detection.LesionDetection} & Dataset lesion AP and FROC for fixed candidates ranked by confidence \\
\end{longtable}
}

By default, probability metrics accept index targets, optional masks and ignored labels. C=1
is binary sigmoid; C\textgreater=2 is mutually exclusive multiclass. Use \texttt{multilabel=True} for independent
channels, as described below. Stateful probability scores average cases; \texttt{Calibration} pools
voxels. All ranks must participate in distributed computation.

\section{Independent probability channels}\label{metrics:independent-probability-channels}

Brier, NLL, ECE, reliability diagrams and their stateful wrappers accept \texttt{multilabel=True}.
Prediction and binary target both have shape \texttt{(B,C,H,W)} or
\texttt{(B,C,D,H,W)}. Channels may overlap; \texttt{from\_logits=True} applies independent sigmoids.
Masks accept \texttt{(B,*spatial)}, \texttt{(B,1,*spatial)} or \texttt{(B,C,*spatial)}. Excluded entries are not
validated. Use masks rather than \texttt{ignore\_index} for independent channels.

For observed voxels \(V_c\) of label \(c\), the per-label scores are:

\[
\mathrm{Brier}_c=\frac{1}{|V_c|}\sum_{i\in V_c}(p_{ic}-y_{ic})^2,\qquad
\mathrm{NLL}_c=-\frac{1}{|V_c|}\sum_{i\in V_c}
\left[y_{ic}\log p_{ic}+(1-y_{ic})\log(1-p_{ic})\right].
\]

Correct endpoint probabilities contribute zero; impossible outcomes have infinite NLL.
Logit input uses stable softplus arithmetic. These are probability quality scores, not
pure calibration measures. Classwise reliability compares positive probability with
positive-label frequency; top-label calibration is rejected for multilabel inputs.

The \texttt{average} option controls within-case reduction. The following definitions also apply
when each label has a different observation mask:

{\small\def\LTcaptype{none} % do not increment counter
\begin{longtable}[]{@{}lll@{}}
\toprule\noalign{}
Option & Output & Reduction \\
\midrule\noalign{}
\endhead
\bottomrule\noalign{}
\endlastfoot
\texttt{"none"} (default) & \texttt{(B,C)} & One score per label; unobserved labels are NaN \\
\texttt{"macro"} & \texttt{(B,1)} & Equal mean of observed label scores \\
\texttt{"micro"} & \texttt{(B,1)} & Pool observed label-voxel pairs before calculating the score \\
\end{longtable}
}

For micro ECE, bin counts, probability sums and positive counts are pooled across labels
\textbf{before} taking absolute calibration gaps. Opposing label errors can cancel; macro ECE
does not permit this cancellation. Fully excluded cases are NaN. Stateful score classes
then average cases; \texttt{Calibration} pools voxels across cases while retaining separate labels.
C++/OpenMP and supported CUDA Triton paths fuse validation and statistics in float64;
other paths use checked PyTorch.

\begin{lstlisting}[language=Python]
import torch
from exacta import functional as EF
from exacta.metrics import Calibration

generator = torch.Generator().manual_seed(5)
logits = torch.randn(2, 3, 8, 8, generator=generator)
labels = torch.rand(2, 3, 8, 8, generator=generator) > 0.7
per_label = EF.brier_score(logits, labels, from_logits=True, multilabel=True)
calibration = Calibration(multilabel=True, from_logits=True, bins=15)
calibration.update(logits, labels)
diagram = calibration.compute()
\end{lstlisting}

\section{Physical area and error}\label{metrics:physical-area-and-error}

The 2D APIs \texttt{predicted\_area}, \texttt{target\_area} and \texttt{area\_error} measure occupied pixel-cell areas.
Their stateful counterparts are \texttt{PredictedArea}, \texttt{TargetArea} and \texttt{AreaError}. With pixel
spacing \((s_H,s_W)\) in millimetres, \(a=s_Hs_W\) and

\[
A_P=a|P|,\qquad A_T=a|T|,\qquad E=A_P-A_T.
\]

The default \texttt{unit="mm2"} returns mm²; \texttt{unit="cm2"} divides by 100. \texttt{absolute=True} on the
error reports \(|E|\). Supply scalar, shared two-axis or per-case \texttt{(B,2)} spacing. Areas require
2D inputs. Mask/ignore policies match physical volume: absent observed classes return zero;
unobserved domains return NaN. Overlapping channels are measured independently.

\texttt{evaluate} and \texttt{MetricCollection} share counts across areas and overlap scores. See the
\href{https://github.com/rodolfo-viana/exacta/wiki/Medical-segmentation-2D}{2D medical-segmentation guide}
for runnable examples and the distinction between slice and patient averages.

\section{Physical volume and error}\label{metrics:physical-volume-and-error}

\texttt{predicted\_volume}, \texttt{target\_volume} and \texttt{volume\_error} take hard 3D
segmentations and spacing in millimetres. Their stateful counterparts are \texttt{PredictedVolume},
\texttt{TargetVolume} and \texttt{VolumeError}. For axis spacings \(s_1,s_2,s_3\), the voxel volume is
\(v=s_1s_2s_3\) mm³. The reported quantities are:

\[
V_P=(\mathrm{TP}+\mathrm{FP})v,\qquad V_T=(\mathrm{TP}+\mathrm{FN})v,\qquad
E=(\mathrm{FP}-\mathrm{FN})v.
\]

\texttt{absolute=True} reports \(|E|\). \texttt{unit="ml"} divides by 1000; the default is \texttt{"mm3"}.
Spacing accepts a scalar, three shared axis spacings or \texttt{(B,3)} per-case spacings.
Two-dimensional areas are rejected. Masking and ignored labels restrict the measured
domain: zero volume means an absent class in an observed domain, while an entirely
unobserved domain returns NaN. These metrics do not infer the volume of unannotated tissue.

\texttt{evaluate} and \texttt{MetricCollection} share one count pass with ordinary overlap metrics:

\begin{lstlisting}[language=Python]
pred = torch.zeros(2, 8, 8, 8, dtype=torch.int64)
target = pred.clone()
target[:, 2:5, 2:5, 2:5] = 1
values = EF.evaluate(pred, target, num_classes=2, include_background=False,
                     spacing=[[1, 1, 2], [0.8, 0.8, 3]],
                     metrics={"dice": {}, "predicted_volume": {"unit": "ml"},
                              "target_volume": {"unit": "ml"},
                              "volume_error": {"unit": "ml", "absolute": True}})
\end{lstlisting}

\section{Voxel PR, AP and AUROC}\label{metrics:voxel-pr-ap-and-auroc}

\texttt{voxel\_ranking} returns per-case/class AP and AUROC arrays plus PR/ROC curves indexed by
\texttt{result.curves{[}case{]}{[}class{]}}. \texttt{average\_precision} and \texttt{auroc} return only the respective
score arrays. C=1 is binary foreground; C\textgreater=2 is categorical one-versus-rest unless
\texttt{multilabel=True} selects independent channels with same-shaped binary targets. The input
and exclusion conventions match the probability APIs above. Categorical targets are indices.

Exact mode (\texttt{bins=None}) sorts scores on the input device and groups equal scores. With
decreasing thresholds, let \(r_j\) and \(q_j\) be recall and precision after adding tie group \(j\):

\[
\mathrm{AP}=\sum_j(r_j-r_{j-1})q_j.
\]

This is the \href{https://scikit-learn.org/stable/modules/generated/sklearn.metrics.average_precision_score.html}{non-interpolated AP convention},
not trapezoidal area under the PR curve. AUROC integrates the ROC curve trapezoidally,
equivalently giving half credit to tied positive-negative score pairs. AP is NaN without
positives and 1 for an observed all-positive target. AUROC is NaN unless both outcomes occur.
Entirely excluded inputs give NaN scores. These are voxel measures, distinct from lesion AP.

Curves start with threshold \texttt{+inf}, precision 1 and recall 0 when positives exist. Thresholds
decrease and inclusion means \texttt{score\ \textgreater{}=\ threshold}. Exact thresholds use the supplied
probabilities; logit input uses raw binary/multilabel logits or categorical log-softmax
scores. Keeping log scores avoids sigmoid/softmax saturation creating artificial ties.
Floating-point normalization can still affect almost-equal categorical scores.

\texttt{bins=256}, for example, quantizes probabilities into 256 equal-width bins; 1 belongs to
the last bin. This is approximate ranking: scores in a bin become tied. Curves use bin
lower boundaries as probability thresholds. More bins reduce quantization but do not
guarantee a monotonic error decrease or a fixed error bound.

\texttt{AveragePrecision} and \texttt{AUROC} average case scores. \texttt{VoxelRanking} instead pools observed
voxels across batches and ranks, so larger images contribute more observations:

\begin{lstlisting}[language=Python]
from exacta.metrics import VoxelRanking

ranking = VoxelRanking(from_logits=True, multilabel=True, bins=256)
ranking.update(logits, labels)
pooled = ranking.compute()  # AP/AUROC have shape (1,C), one pooled dataset
precision = pooled.curves[0][0].precision
recall = pooled.curves[0][0].recall
ranking.reset()
\end{lstlisting}

Exact streaming retains selected scores and labels and gathers each class on every rank.
Binned streaming retains \texttt{2*C*bins} int64 counts and sums those counts across ranks;
temporary update buffers depend on the current batch. Every rank, including empty ranks,
must call \texttt{compute} with identical configuration. Set the current CUDA device for NCCL.
Repeated \texttt{compute} does not add observations or mutate local state.

% Adapted from docs/wiki/Losses.md; edit this file, not the wiki page, for the manual.
\chapter{Losses}\label{ch:losses}

Losses are PyTorch modules accepting floating-point logits \texttt{(B,\ C,\ *spatial)}. Most accept integer
targets \texttt{(B,\ *spatial)} and \texttt{target\_format="one-hot"} for same-shaped masks or soft labels;
individual objectives document their supported target formats below.
Softmax targets should sum to one over classes; sigmoid targets should lie in \texttt{{[}0,\ 1{]}}.
Exacta does not silently normalize or clamp target probabilities.

Let \texttt{p} be activated predictions, \texttt{g} the target and \texttt{s} smoothing. Core overlap and focal losses are:

{\small\def\LTcaptype{none} % do not increment counter
\begin{longtable}[]{@{}
  >{\raggedright\arraybackslash}p{(\linewidth - 2\tabcolsep) * \real{0.2700}}
  >{\raggedright\arraybackslash}p{(\linewidth - 2\tabcolsep) * \real{0.7300}}@{}}
\toprule\noalign{}
\begin{minipage}[b]{\linewidth}\raggedright
Loss
\end{minipage} & \begin{minipage}[b]{\linewidth}\raggedright
Definition and reduction
\end{minipage} \\
\midrule\noalign{}
\endhead
\bottomrule\noalign{}
\endlastfoot
\texttt{DiceLoss} & \texttt{1\ -\ (2\ sum(p*g)+s)/(sum(p)+sum(g)+s)}, averaged across cases and selected classes \\
\texttt{DiceCELoss} & \texttt{lambda\_dice\ *\ DiceLoss\ +\ lambda\_ce\ *\ CE}; sigmoid uses BCE \\
\texttt{DiceFocalLoss} & \texttt{lambda\_dice\ *\ DiceLoss\ +\ lambda\_focal\ *\ FocalLoss} \\
\texttt{GeneralizedDiceCELoss} & \texttt{lambda\_dice\ *\ GeneralizedDiceLoss\ +\ lambda\_ce\ *\ CE}; sigmoid uses BCE \\
\texttt{TverskyLoss} & \texttt{1\ -\ (TP+s)/(TP+alpha*FP+beta*FN+s)}, averaged across cases and selected classes \\
\texttt{GeneralizedDiceLoss} & Ratio of class-weighted overlap sums; class weights derive from target volumes \\
\texttt{FocalLoss} & Focal CE/BCE with exponent \texttt{gamma}; softmax sums classes then averages voxels, sigmoid averages all entries \\
\end{longtable}
}

Overlap losses support \texttt{activation="softmax"}, \texttt{"sigmoid"} or \texttt{"none"}. Dice-CE and focal
require logits and use softmax or sigmoid. \texttt{batch\_dice=True} pools overlap statistics across
the local batch before division. The explicit global option is documented below.
\texttt{include\_background=False} excludes channel zero from overlap only. CE/BCE and focal include
all channels.

\section{Overlap equations}\label{losses:overlap-equations}

For one case and selected class, let \(p_i\) be the activated prediction and \(g_i\) the target
at included voxel \(i\). With smoothing \(s\), the soft overlap statistics are:

\[
\mathrm{TP}=\sum_i p_i g_i,\qquad
\mathrm{FP}=\sum_i p_i(1-g_i),\qquad
\mathrm{FN}=\sum_i (1-p_i)g_i.
\]

Dice and Tversky use these ratios before averaging over cases and selected classes:

\[
\mathcal{L}_{\mathrm{Dice}}
=1-\frac{2\sum_i p_i g_i+s}{\sum_i p_i+\sum_i g_i+s},\qquad
\mathcal{L}_{\mathrm{Tversky}}
=1-\frac{\mathrm{TP}+s}{\mathrm{TP}+\alpha\mathrm{FP}+\beta\mathrm{FN}+s}.
\]

Generalized Dice instead takes a single class-weighted ratio per case, where \(c\) ranges
over selected classes:

\[
\mathcal{L}_{\mathrm{GDice}}
=1-\frac{2\sum_c w_c\sum_i p_{ic}g_{ic}+s}
{\sum_c w_c\sum_i(p_{ic}+g_{ic})+s}.
\]

For present target classes, \texttt{weight\_type="square"} sets \(w_c=(\sum_i g_{ic})^{-2}\);
\texttt{"simple"} uses the inverse volume and \texttt{"uniform"} uses one. The absent-class rules below
handle zero volumes. \texttt{batch\_dice=True} includes the batch axis in the sums before taking
the ratio. Smoothing is added once per ratio, after weighting for generalized Dice.

\section{Important parameters}\label{losses:important-parameters}

The following options describe the core objectives. Squared Dice, Lovasz-hinge and
Wasserstein Dice are described at the end of this page.

Definitions differ across libraries, so select parameters intentionally:

\begin{itemize}
\tightlist
\item
  Dice/Tversky smoothing defaults to \texttt{1e-5}. Tversky's default \texttt{alpha=beta=0.5} equals Dice
  with twice the smoothing constant; alpha weights false positives and beta false negatives.
\item
  \texttt{zero\_absent=True} (Dice, Tversky and focal Tversky, since 0.1.0) gives zero loss and
  gradient to classes whose target is empty within the reduction (case, or batch with
  \texttt{batch\_dice=True}), however much is predicted; they still count in the mean. This reproduces
  segmentation\_models\_pytorch's Dice and Tversky losses with \texttt{batch\_dice=True} and \texttt{smooth=0}.
  By default, as in MONAI, an empty target still penalizes predicted probability. It runs on the
  PyTorch backend.
\item
  Generalized Dice supports \texttt{weight\_type="square"}, \texttt{"simple"} and \texttt{"uniform"}. Smoothing
  is added after class weighting. Absent classes receive the largest present-class weight.
  With inverse-volume weighting, \texttt{empty\_target="zero"} gives zero for entirely absent selected
  targets; \texttt{"uniform"} uses unit weights there. Positive smoothing is needed for a useful
  gradient on an entirely empty target. This differs from an entirely ignored case.
\item
  Focal uses \texttt{gamma=2} and \texttt{alpha=None} by default. Scalar alpha weights sigmoid positives by
  alpha and negatives by \texttt{1-alpha}; softmax assigns \texttt{1-alpha} to class zero and alpha to other
  classes. \texttt{gamma=0,\ alpha=None} reproduces CE/BCE. Softmax reduction differs from MONAI's mean
  over channels by a factor of the channel count.
\end{itemize}

\section{Backends and derivatives}\label{losses:backends-and-derivatives}

\texttt{backend="auto"} selects fused Triton for supported CUDA inputs and PyTorch otherwise.
Half/bfloat16 inputs compute in float32; float64 inputs retain their precision. Fused losses
support first derivatives. Use \texttt{backend="torch"} for higher-order derivatives; \texttt{auto} also
selects PyTorch when targets require gradients. Full-graph \texttt{torch.compile} is tested for
local losses. Distributed losses are tested in eager execution; compiled collectives are
not part of the current compatibility guarantee.

\section{Ignored labels and voxel masks}\label{losses:ignored-labels-and-voxel-masks}

Dice, Tversky, generalized Dice, focal and their regional combinations accept constructor
argument \texttt{ignore\_index} for index targets and keyword-only \texttt{mask=} in forward. Focal Tversky,
top-k CE and Lovasz-Softmax also support exclusions. Geometry/topology losses have different
requirements, described below.
Masks must be boolean, on the input device, with shape \texttt{(B,\ *spatial)} or \texttt{(B,\ 1,\ *spatial)}.
True includes a voxel; False excludes it from every class. Both exclusions combine.

\begin{lstlisting}[language=Python]
import torch
from exacta.losses import DiceCELoss

generator = torch.Generator().manual_seed(0)
logits = torch.randn(2, 3, 4, 4, 4, generator=generator, requires_grad=True)
target = torch.randint(3, (2, 4, 4, 4), generator=generator)
target[:, 0] = -100
mask = torch.ones_like(target, dtype=torch.bool)
mask[:, :, 0] = False
loss = DiceCELoss(ignore_index=-100)(logits, target, mask=mask)
loss.backward()
assert torch.count_nonzero(logits.grad[:, :, 0]) == 0
\end{lstlisting}

CE/focal normalize over included voxels (and channels for sigmoid). Overlap averages exclude
entirely ignored cases; a fully ignored batch returns zero loss and gradients even with zero
smoothing. Annotated background voxels remain valid observations. Ignored invalid labels and
nonfinite values are excluded before activation/reduction. Valid out-of-range labels still
raise eagerly; compiled Triton returns NaN. Soft/one-hot targets use \texttt{mask}, not \texttt{ignore\_index}.

The GPU kernels read exclusion information directly and count valid voxels in their existing
passes. They do not create full-volume masked probability or one-hot buffers.

\section{Loss combinations and CE weights}\label{losses:loss-combinations-and-ce-weights}

\texttt{DiceFocalLoss} adds Dice and focal terms; \texttt{GeneralizedDiceCELoss} adds generalized Dice and
CE/BCE. Their GPU implementations share one activation/target pass forward and one backward.
The corresponding native PyTorch compositions provide a reference and higher-order gradients.
Coefficients must be finite and nonnegative, with at least one positive term. A coefficient
of zero skips that term, including otherwise undefined zero-smoothing overlap ratios.

\begin{lstlisting}[language=Python]
from exacta.losses import DiceFocalLoss, GeneralizedDiceCELoss

dice_focal = DiceFocalLoss(include_background=False, gamma=2, alpha=0.25,
                          lambda_dice=1, lambda_focal=0.5)
gdl_ce = GeneralizedDiceCELoss(include_background=False, empty_target="uniform",
                              class_weight=[0.2, 1.0, 2.0])
\end{lstlisting}

\texttt{DiceCELoss} and \texttt{GeneralizedDiceCELoss} accept \texttt{class\_weight}, one nonnegative finite value
per input channel, stored as a module buffer. These weights apply only to CE/BCE. Normalization
matches \href{https://docs.pytorch.org/docs/2.14/generated/torch.nn.CrossEntropyLoss.html}{PyTorch's cross-entropy definition}
and depends on target format:

{\small\def\LTcaptype{none} % do not increment counter
\begin{longtable}[]{@{}
  >{\raggedright\arraybackslash}p{(\linewidth - 4\tabcolsep) * \real{0.3000}}
  >{\raggedright\arraybackslash}p{(\linewidth - 4\tabcolsep) * \real{0.4500}}
  >{\raggedright\arraybackslash}p{(\linewidth - 4\tabcolsep) * \real{0.2500}}@{}}
\toprule\noalign{}
\begin{minipage}[b]{\linewidth}\raggedright
Pointwise term
\end{minipage} & \begin{minipage}[b]{\linewidth}\raggedright
Weighted numerator
\end{minipage} & \begin{minipage}[b]{\linewidth}\raggedright
Denominator
\end{minipage} \\
\midrule\noalign{}
\endhead
\bottomrule\noalign{}
\endlastfoot
Softmax, index target & Sum of \texttt{-w{[}y{]}\ log\ p{[}y{]}} over included voxels & Sum of included \texttt{w{[}y{]}} \\
Softmax, soft/one-hot target & Sum of \texttt{-w{[}c{]}\ g{[}c{]}\ log\ p{[}c{]}} & Number of included voxels \\
Sigmoid BCE & Class weight multiplies both positive and negative terms & Included voxels × channels \\
\end{longtable}
}

Zero total target weight returns zero CE and zero CE gradient. Exact one-hot and index labels
therefore produce different weighted CE scales, as in PyTorch. \texttt{include\_background=False}
still affects only the overlap term. Use focal's \texttt{alpha} for its documented positive/negative
or background/foreground weighting.

\section{Deep supervision}\label{losses:deep-supervision}

\texttt{DeepSupervisionLoss} applies one scalar loss to explicitly paired resolutions. Weights are
used as supplied, without normalization; zero-weight levels are not evaluated and receive
no gradient. Targets and optional masks must already match each prediction's spatial shape.
The wrapper does not resize labels or choose downsampling conventions.

\begin{lstlisting}[language=Python]
from exacta.losses import DeepSupervisionLoss, DiceCELoss

loss_fn = DeepSupervisionLoss(DiceCELoss(ignore_index=-100), weights=[1.0, 0.5, 0.25])
# Sequences of three tensors, ordered consistently from fine to coarse:
loss = loss_fn(predictions, targets, masks=voxel_masks)
\end{lstlisting}

\section{Global distributed losses}\label{losses:global-distributed-losses}

Set \texttt{distributed=True} to pool statistics over the default \texttt{torch.distributed} process group.
Overlap losses and their combinations also require \texttt{batch\_dice=True}. This computes the
objective for the concatenated global batch, including generalized-Dice class weights and
weighted CE denominators. It differs from averaging independently computed rank-local Dice.

\begin{lstlisting}[language=Python]
# After initializing the process group and wrapping the model in ordinary DDP:
loss_fn = GeneralizedDiceCELoss(batch_dice=True, distributed=True,
                               include_background=False, ignore_index=-100)
loss = loss_fn(ddp_model(images), labels, mask=annotated)
loss.backward()
\end{lstlisting}

Every rank must use identical loss configuration and backend, and call forward and backward
in the same order, including ranks with zero cases or entirely ignored data. An entirely
ignored or empty global batch returns zero. Do not multiply the loss by local batch size or
divide by world size: the differentiable statistics sum and DDP's gradient averaging together
produce concatenated-batch \textbf{parameter gradients}; see
\href{https://docs.pytorch.org/docs/2.14/generated/torch.nn.parallel.DistributedDataParallel.html}{DDP's gradient-averaging convention}.
Logit/target gradients before DDP averaging
carry a world-size factor; custom gradient communication hooks need their own scaling analysis.
Gradient accumulation across microbatches averages microbatch objectives; it does not pool
their nonlinear Dice ratios into one larger-batch ratio.

Only small statistics communicate. NCCL keeps these on the rank's CUDA device. Gloo uses CPU
buffers, also allowing two CUDA workers on a single GPU for correctness tests. Local tests
cover CPU and fused CUDA with two Gloo ranks, uneven batches, empty ranks, ignored ranks and
all-ignored batches. \textbf{The multi-GPU NCCL path is experimental and unverified in 0.1.2:}
the two-GPU tests have not been run on suitable hardware. Gloo checks do not establish NCCL
correctness or multi-GPU training throughput. Loss configuration and valid input shapes
must agree across ranks; as with ordinary distributed training, a rank-local error can interrupt its peers.

\section{Boundary, topology and hard-example objectives}\label{losses:boundary-topology-and-hard-example-objectives}

Additional objectives include \texttt{FocalTverskyLoss}, \texttt{TopKCELoss}, \texttt{DiceTopKCELoss}, \texttt{LovaszSoftmaxLoss},
\texttt{BoundaryLoss}, \texttt{HausdorffDTLoss}, \texttt{SoftClDiceLoss} and \texttt{SkeletonRecallLoss}. Their formulas, reductions, defaults,
reference conventions and CPU/GPU paths are documented in the
\href{https://github.com/rodolfo-viana/exacta/wiki/Advanced-segmentation}{advanced segmentation guide}.

Focal Tversky supports the global distributed reduction above. Top-k/Lovasz selection and
geometry/topology losses operate on each rank's local batch. Geometry masks affect reduction
only; complete target geometry is required. Geometry/topology losses have no \texttt{ignore\_index},
and soft-clDice has no mask argument. Distance-transform and skeleton objectives should be
combined with regional supervision according to the task.

\section{Soft labels}\label{losses:soft-labels}

With soft targets \(g_i\in(0,1)\), the product overlap \(\sum_i p_i g_i\) makes Dice and Tversky
losses minimal away from \(p=g\): they push predictions towards 0 and 1. Label smoothing,
distillation and averaged annotations therefore need another form. \texttt{soft\_label=True} in
\texttt{DiceLoss}, \texttt{TverskyLoss}, \texttt{FocalTverskyLoss}, \texttt{GeneralizedDiceLoss}, \texttt{DiceCELoss},
\texttt{DiceFocalLoss} and \texttt{GeneralizedDiceCELoss} uses the semimetric statistics of
\href{https://doi.org/10.1007/978-3-031-43898-1_46}{Wang et al.~(MICCAI 2023)}:

\[
\mathrm{TP}=\tfrac{1}{2}\Big(\sum_i p_i+\sum_i g_i-\sum_i\lvert p_i-g_i\rvert\Big),\qquad
\mathrm{FP}=\sum_i p_i-\mathrm{TP},\qquad \mathrm{FN}=\sum_i g_i-\mathrm{TP}.
\]

For binary targets \(\lvert p_i-g_i\rvert=p_i+g_i-2p_ig_i\), so values equal the ordinary losses
exactly, and so do gradients except at voxels where \(p_i=g_i\in\{0,1\}\), where the zero
subgradient of \(\lvert p_i-g_i\rvert\) applies; through softmax or sigmoid, whose derivative
vanishes there, logit gradients agree. For soft targets every ratio is maximal only at \(p=g\) (for Tversky,
when both weights are positive). It matches MONAI's \texttt{soft\_label=True}. The option acts on
one-hot targets and runs on the PyTorch backend; index targets are hard and keep the fused
path. CE terms are
unchanged: cross-entropy is already minimal at \(p=g\). \texttt{squared\_pred=True} is too
(\(\sum p^2+\sum g^2-2\sum pg=\sum(p-g)^2\)), so the two options are exclusive.

\begin{lstlisting}[language=Python]
from exacta.losses import DiceCELoss

loss_fn = DiceCELoss(target_format="one-hot", soft_label=True)
loss = loss_fn(logits, smoothed_one_hot_targets)
\end{lstlisting}

\section{Squared Dice denominator}\label{losses:squared-dice-denominator}

\texttt{squared\_pred=True} is available in \texttt{DiceLoss}, \texttt{DiceCELoss}, \texttt{DiceFocalLoss}
and \texttt{DiceTopKCELoss}. Both prediction and target terms in the denominator are squared:

\[
\mathcal{L}_{\mathrm{squared\ Dice}}
=1-\frac{2\sum_i p_i g_i+s}{\sum_i p_i^2+\sum_i g_i^2+s}.
\]

The numerator retains the product overlap, including for soft targets. With binary targets,
\(g_i^2=g_i\). The ordinary denominator remains the default. Smoothing, background exclusion,
masking and case/batch reduction follow ordinary Dice. The option is named consistently with
\href{https://monai.readthedocs.io/en/stable/losses.html\#diceloss}{MONAI's squared-prediction option};
the equation above defines Exacta's soft-target convention explicitly.

\begin{lstlisting}[language=Python]
from exacta.losses import DiceCELoss

loss_fn = DiceCELoss(squared_pred=True, include_background=False)
\end{lstlisting}

PyTorch supports differentiation through predictions and soft targets. The fused CUDA path
includes the squared denominator in both forward statistics and analytic backward, and
supports global distributed Dice/CE reductions. Fused higher-order derivatives remain
unsupported; use \texttt{backend="torch"} for those. Dice-top-k still selects hard examples locally.

\section{Binary and multilabel Lovasz-hinge}\label{losses:binary-and-multilabel-lovasz-hinge}

\texttt{LovaszHingeLoss} follows the binary real-valued-score formulation in the
\href{https://github.com/bermanmaxim/LovaszSoftmax/blob/master/pytorch/lovasz_losses.py}{authors' implementation}.
Given binary labels \(y_i\) and raw logits \(z_i\), margin errors are
\(e_i=1-z_i(2y_i-1)\). Sort the errors decreasingly, and let \(\Delta J_j\) be successive
increments of the Jaccard error for that ordering. The objective is:

\[
\mathcal{L}_{\mathrm{Lovasz\ hinge}}=\sum_j\max(e_{(j)},0)\,\Delta J_j.
\]

Binary input has shape \texttt{(B,1,*spatial)} with target \texttt{(B,*spatial)}. Independent overlapping
labels use \texttt{multilabel=True} and a same-shaped binary target. No sigmoid is applied and
soft targets are rejected. Masks can be shared or per-label; \texttt{ignore\_index} applies only
to integer binary targets. Excluded invalid logits/targets do not affect the result.

\texttt{per\_image=True} (default) averages observed case/label losses. False pools the local batch
per label before averaging labels. An absent positive class still penalizes false positives;
unobserved labels are omitted, and entirely excluded inputs give differentiable zero.
Stable native PyTorch sorting uses flattened voxel order for ties. The subgradient at a
zero hinge margin is zero. This is a local objective under DDP, without global sorting.

\begin{lstlisting}[language=Python]
from exacta.losses import LovaszHingeLoss

binary_loss = LovaszHingeLoss(ignore_index=-1)
region_loss = LovaszHingeLoss(multilabel=True, per_image=True)
\end{lstlisting}

These losses support both 2D and 3D medical segmentation. They are alternative
objectives with different gradients, not a claim of improved accuracy for every task.

\section{Generalized Wasserstein Dice}\label{losses:generalized-wasserstein-dice}

\texttt{GeneralizedWassersteinDiceLoss} adds explicit costs for confusing mutually exclusive labels.
It follows the crisp-target formulation of \href{https://arxiv.org/abs/1707.00478}{Fidon et al.~(2017)},
with the alternative \texttt{GDL} weighting described by
\href{https://monai.readthedocs.io/en/stable/losses.html\#generalizedwassersteindiceloss}{MONAI}.
It supports 2D images and 3D volumes. Costs describe relationships between labels; voxel
spacing and physical boundary distances do not enter this objective.

Let \(M\) be the supplied class-cost matrix divided by its positive maximum, \(y_i\) the integer
target at observed voxel \(i\), and \(p_{ic}\) the softmax probability for class \(c\). The expected
cost against that crisp target is

\[
w_i=\sum_c M_{y_i,c}\,p_{ic}.
\]

For each case, define \(N_c=\sum_i\mathbf{1}[y_i=c]\) over the observed domain and
\(T=\sum_i\alpha_{y_i}(1-w_i)\). The two modes use different weights and error sums:

{\small\def\LTcaptype{none} % do not increment counter
\begin{longtable}[]{@{}lll@{}}
\toprule\noalign{}
\texttt{weighting\_mode} & Class weights \(\alpha_c\) & Error \(E\) \\
\midrule\noalign{}
\endhead
\bottomrule\noalign{}
\endlastfoot
\texttt{"default"} & \(\alpha_0=0\); \(\alpha_c=1\) for \(c>0\) & \(\sum_i w_i\) \\
\texttt{"GDL"} & \(1/(N_c+1)\), including label zero & \(\sum_i\alpha_{y_i}w_i\) \\
\end{longtable}
}

Both return the following loss, with the same smoothing constant \(s\) in both terms:

\[
\mathcal{L}=1-\frac{2T+s}{2T+E+s}.
\]

The constructor requires a finite, nonnegative, symmetric square matrix with zero diagonal
and a positive maximum. \texttt{default} additionally requires maximal costs between background
(label zero) and every foreground class, matching the assumption behind the simplified
foreground formula. \texttt{GDL} permits other background costs. Triangle inequalities are not
checked; a zero off-diagonal cost makes that label confusion cost-free. Specify these costs
from the task's label relationships, and assess any tuning using training/validation data.

The following synthetic example gives confusion between foreground labels a smaller cost
than confusion with background. The same loss accepts \texttt{(B,C,H,W)} or \texttt{(B,C,D,H,W)} logits.

\begin{lstlisting}[language=Python]
import torch
from exacta.losses import GeneralizedWassersteinDiceLoss

device = torch.device("cuda" if torch.cuda.is_available() else "cpu")
rng = torch.Generator().manual_seed(879)
logits = torch.randn((2, 3, 16, 20), generator=rng).to(device).requires_grad_()
target = torch.randint(3, (2, 16, 20), generator=rng).to(device)
target[0, :2] = -1

loss_fn = GeneralizedWassersteinDiceLoss(
    [[0, 1, 1], [1, 0, 0.25], [1, 0.25, 0]],
    weighting_mode="default", ignore_index=-1,
).to(device=device, dtype=logits.dtype)
loss = loss_fn(logits, target)
loss.backward()
assert torch.isfinite(loss) and torch.isfinite(logits.grad).all()
\end{lstlisting}

The default \texttt{reduction="mean"} averages observed cases equally. \texttt{"sum"} sums case losses;
\texttt{"none"} returns \texttt{(B,)}. Fully excluded cases and empty inputs return differentiable zero,
and are omitted from the mean. A zero denominator with \texttt{smooth=0} also returns zero.
With default weights and \texttt{smooth=0}, an observed background-only case with any nonzero error
has loss one and no useful foreground Dice gradient; positive smoothing or the \texttt{GDL} mode
changes that behavior. This differs from an entirely unobserved case.

This API takes integer targets and logits; it has no sigmoid/multilabel, soft-target or
cross-rank pooling option. It supports shared boolean masks, ignored indices, half/bfloat16
computation in float32, float64, higher derivatives and full-graph \texttt{torch.compile} on CPU/CUDA.
Eager calls reject out-of-range observed labels. Compiled calls return NaN for invalid labels
without unsafe indexing. Excluded labels/nonfinite logits are sanitized before activation.

The implementation uses native PyTorch gather/reduction kernels and an integer class histogram,
without expanding targets to one-hot. Move the loss to the logits' device/dtype once and
benchmark \texttt{torch.compile(loss\_fn)} with forward \textbf{and backward} on your workload. The cost
matrix is a detached normalized buffer, included in \texttt{state\_dict}; it is not a learnable parameter.

\chapter{Reproducing the definitions of other tools}\label{ch:tools}

The surface scores of Chapter~\ref{ch:conventions} have arguments for every choice in which
common tools differ. Matching a tool means passing the arguments in
Table~\ref{tab:presets}. The example driver (Chapter~\ref{ch:driver}) keeps them in a
dictionary named \texttt{PRESETS}; the referee tests (Chapter~\ref{ch:tests}) compare each
preset with the tool itself.

\begin{table}[ht]
\centering
\footnotesize
\begin{tabular}{@{}>{\raggedright\arraybackslash}p{0.18\linewidth}%
  >{\raggedright\arraybackslash}p{0.22\linewidth}%
  >{\raggedright\arraybackslash}p{0.22\linewidth}%
  >{\raggedright\arraybackslash}p{0.28\linewidth}@{}}
\toprule
Tool & HD95 & ASSD & NSD \\
\midrule
MONAI &
\texttt{surface="voxels"}, \texttt{combine="max"} &
\texttt{surface="voxels"}, \texttt{combine="pooled"} &
\texttt{surface="voxels"}; with \texttt{use\_subvoxels=True} use \texttt{surface="surfels"} \\
MedPy 0.5.2 &
\texttt{surface="voxels"}, \texttt{combine="pooled"} &
\texttt{surface="voxels"}, \texttt{combine="pooled"} &
(not provided) \\
MedPy 0.4 &
as 0.5.2 &
\texttt{combine="mean"} &
(not provided) \\
DeepMind \texttt{surface-distance} &
\texttt{surface="surfels"}, \texttt{combine="max"} &
\texttt{surface="surfels"}; \texttt{combine="mean"} is the mean of the two directed means it returns &
\texttt{surface="surfels"}, \texttt{tie\_rtol=0.0} \\
MeshMetrics 1.x &
\texttt{surface="mesh"}, \texttt{combine="max"} &
\texttt{surface="mesh"}, \texttt{combine="pooled"} (\texttt{assd}) or \texttt{"mean"} (\texttt{masd}) &
\texttt{surface="mesh"} \\
\bottomrule
\end{tabular}
\caption{Arguments of \texttt{exacta.functional.hd95}, \texttt{assd} and \texttt{nsd} that
reproduce the definitions of other tools. Everything not listed keeps its default.}
\label{tab:presets}
\end{table}

For example, MedPy's pooled HD95 is
\begin{lstlisting}[language=Python]
import exacta.functional as EF

medpy = dict(surface="voxels", combine="pooled")
scores = EF.hd95(prediction, target, spacing=spacing, num_classes=2,
                 include_background=False, **medpy)
\end{lstlisting}
and DeepMind's surface Dice at a tolerance of 1.2 is
\texttt{EF.nsd(prediction, target, spacing=spacing, tolerance=1.2, surface="surfels", tie\_rtol=0.0, \ldots)}.

\section{What still differs}\label{tools:differences}

\paragraph{Ties with the NSD tolerance.} A surface element whose distance equals the tolerance
in exact arithmetic can fall on either side of it after rounding. Exacta's default margin
$\rho=10^{-6}$ counts such ties as within the tolerance (Chapter~\ref{ch:conventions}).
\texttt{tie\_rtol=0} compares strictly in binary64, as SciPy-based tools (MedPy and DeepMind's
\texttt{surface-distance}) do; each distance is then the exact floating-point minimum of the
expression that SciPy's \texttt{distance\_transform\_edt} evaluates, so it lands on the same side
of the tolerance. SciPy itself returns a value one unit in the last place above that minimum at
a few voxels per million, and a tie at such a voxel can still differ. MONAI casts distances to
float32 before comparing them with the tolerance, which counts exact ties as within, as the
default margin does, and also distances within float32 resolution of the tolerance.

\paragraph{MeshMetrics.} Exacta meshes each mask in the box that MeshMetrics 1.2 to 1.4.1
use for NumPy masks (the bounding box of both masks, starting one voxel before it), with VTK's
discrete marching cubes (discrete flying edges in 2-D). In 3-D, HD, HD95 and ASSD agree with
MeshMetrics to about $10^{-13}$ in the unit of the spacing. In 2-D, MeshMetrics measures
distances to an extruded contour that VTK stores in float32, so values differ by about
$10^{-7}$ in relative terms. NSD agrees except at ties, triangle centres whose distance to the
other mesh equals the tolerance in exact arithmetic. In the tests, MeshMetrics' NSD lies
between Exacta's values with every tie outside and with every tie within the tolerance.
Bit-for-bit agreement at ties was observed only with exact coordinates (for example spacing 1
or 0.5), NumPy masks, x86-64 and tolerances 1 and 2. The default margin counts every tie as
within, whichever tool and platform, and is the setting to use when a result must not depend on
either.

\paragraph{Versions.} Definitions change between releases of the tools. MedPy's ASSD was the
mean of the two directed means in 0.4 and is the pooled mean in 0.5.2. The 0.2 releases of
MeshMetrics used surface nets, and releases 1.0 to 1.4.1 use discrete marching cubes. Record the
version of every tool together with the Exacta arguments when reporting a comparison.

\section{Checking the presets}\label{tools:checking}

The referee subset (Chapter~\ref{ch:tests}) runs these comparisons when the tools are installed:
MONAI and MedPy in \texttt{test\_surface.py} and \texttt{test\_volume.py}; DeepMind's lookup
tables bit for bit, and HD, HD95, ASSD and NSD with surfels, in \texttt{test\_surfels.py};
MeshMetrics and VTK's marching cubes in \texttt{test\_mesh.py}. A test whose tool is missing is
skipped, and the summary line of pytest says how many were skipped.

% Adapted from docs/wiki/Performance.md; edit this file, not the wiki page, for the manual.
\chapter{Performance and backends}\label{ch:performance}

Exacta uses compiled C++ CPU kernels for expensive metric passes and Triton for supported
CUDA paths. Losses use native PyTorch/autograd and, for supported CUDA objectives,
fused Triton forward/backward.
Backend choice should be supported by measurements for the shapes and dtypes used in training.

Fused losses recompute activations in backward and retain inputs plus small statistics, avoiding
full-volume saved probability and one-hot intermediates. Mesh distances currently run on CPU,
including when input arrays reside on CUDA. Cross-device conversion costs matter in a pipeline.

\section{Reproducing measurements}\label{performance:reproducing-measurements}

The benchmark scripts used for these published measurements remain available in the
\texttt{v0.0.1a3} source snapshot. Inputs are seeded, compilation is warmed up, and GPU steps
synchronize around timing. Clone that snapshot before running the commands below:

\begin{lstlisting}[language=bash]
git clone --branch v0.0.1a3 https://github.com/rodolfo-viana/exacta.git exacta-benchmarks
cd exacta-benchmarks
python -m pip install -e ".[test]"
python benchmarks/bench_losses.py --classes 4 --sizes 128 --dtypes float32 \
  --repeat 20 --warmup 5 --seed 0 --threads 16
python benchmarks/bench_losses.py --device cpu --classes 4 --sizes 64 --dtypes float32 \
  --backends torch compiled --repeat 20 --warmup 5 --seed 0 --threads 16
python benchmarks/bench_metrics.py --sizes 128 256 --only exacta
\end{lstlisting}

Local generalized-Dice measurements for version 0.0.1a2 use float32 softmax inputs, PyTorch 2.14.1,
CUDA 13.0, an RTX 5070 and a Ryzen 9 5900XT with 16 CPU threads. Medians of 20 steps after five
warm-up steps are:

{\small\def\LTcaptype{none} % do not increment counter
\begin{longtable}[]{@{}lrrr@{}}
\toprule\noalign{}
Device and shape & Eager PyTorch & Compiled PyTorch & Fused Triton \\
\midrule\noalign{}
\endhead
\bottomrule\noalign{}
\endlastfoot
GPU, \texttt{2×4×128³} & 4.184 ms & 3.521 ms & 0.938 ms \\
CPU, \texttt{2×4×64³} & 10.541 ms & 4.985 ms & Not applicable \\
\end{longtable}
}

GPU additional peak allocations were 256 MiB, 128 MiB and under 0.1 MiB respectively, measured
beyond inputs and the logits' gradient. Scratch allocations can fit below that baseline.
These are synthetic local measurements, not guaranteed speedups on other hardware or complete
training pipelines. Record code version, shapes, classes, dtype, threads and environment.

\section{Masked losses}\label{performance:masked-losses}

The loss benchmark supports \texttt{-\/-ignore-fraction} and \texttt{-\/-mask-fraction}. At code revision
\href{https://github.com/rodolfo-viana/exacta/commit/d3ba67fbe8b24ec6b65bdc1f1fa3f813fcb6ad9c}{\texttt{d3ba67f}},
independently excluding 25\% of labels and 20\% through a boolean mask retained 60.01\% of voxels
on the GPU input. On the same RTX 5070 with float32 \texttt{2×4×128³} logits, seed 0, five warm-up
steps and 20 measured forward/backward steps, median times were:

{\small\def\LTcaptype{none} % do not increment counter
\begin{longtable}[]{@{}lrrr@{}}
\toprule\noalign{}
Loss & Eager PyTorch & Compiled PyTorch & Fused Triton \\
\midrule\noalign{}
\endhead
\bottomrule\noalign{}
\endlastfoot
Dice & 5.741 ms & 2.479 ms & 1.030 ms \\
Dice + CE & 8.271 ms & 3.719 ms & 1.071 ms \\
Tversky & 8.757 ms & 2.558 ms & 1.068 ms \\
Focal & 9.869 ms & 2.605 ms & 0.728 ms \\
Generalized Dice & 5.760 ms & 2.620 ms & 1.158 ms \\
\end{longtable}
}

On CPU, masked Dice + CE with \texttt{2×4×64³} logits took 23.432 ms eager and 11.087 ms compiled,
using 16 threads. The \href{https://github.com/rodolfo-viana/exacta/blob/v0.0.1a3/benchmarks/results/masking-2026-10-24.md}{complete measurement record}
includes unmasked runs, allocation measurements, environment versions, commands and limitations.

\section{Compiled CPU losses}\label{performance:compiled-cpu-losses}

For repeated CPU training, compiling native PyTorch operations can reduce allocations and
dispatch overhead. Compilation has startup cost and should be excluded from steady-state
comparisons while still being considered for short runs:

\begin{lstlisting}[language=Python]
import torch
from exacta.losses import DiceCELoss

loss_fn = torch.compile(DiceCELoss(backend="torch"), fullgraph=True)
\end{lstlisting}

\section{Shared metrics and composed losses}\label{performance:shared-metrics-and-composed-losses}

At revision \href{https://github.com/rodolfo-viana/exacta/commit/12af1d4762ef03b5724b2603ca26f80051bf9f75}{\texttt{12af1d4}},
float32 Dice--focal on the same RTX 5070, with \texttt{2×4×128³} logits, took 0.861 ms for combined
Triton forward/backward, 1.622 ms for separate fused Dice and focal calls, and 3.390 ms for
compiled native Torch. Additional GPU peak allocation was 0.0, 64.0 and 256.0 MiB respectively,
using the baseline definition above. CPU native Torch also benefited from compilation.

Shared evaluation of seven count metrics plus HD95, ASSD and NSD on two synthetic \texttt{128³}
cuboid cases took 6.714 ms versus 27.696 ms for separate CPU calls, and 5.659 ms versus
15.260 ms for separate CUDA calls. Spacing was anisotropic and background was excluded.
All figures are medians of 20 repetitions after five warm-ups, using 16 threads.

The \href{https://github.com/rodolfo-viana/exacta/blob/v0.0.1a3/benchmarks/results/composition-collection-2026-10-24.md}{measurement record}
contains the full environment, code revision, CPU results, bfloat16 runs, masked class-weighted
CE combinations, exact commands and limitations. The scripts are \texttt{bench\_losses.py} and
\texttt{bench\_collection.py}. Distributed correctness tests do not establish multi-GPU throughput;
NCCL benchmarks require additional hardware.

Avoid running performance measurements alongside other workloads. Measure sparse and dense
targets, different class counts and representative MRI sizes before choosing an implementation.

\section{Probability, ranking and loss measurements}\label{performance:probability-ranking-and-loss-measurements}

The source includes \texttt{bench\_probability.py} and \texttt{bench\_completion.py}. Measurements
on 2026-10-26 used Python 3.14.4, PyTorch 2.14.1, Ryzen 9 5900XT (two threads) and RTX 5070,
float32 inputs of shape \texttt{2×3×64³}, three warmups and ten synchronized repetitions.
The tables report medians on synthetic inputs; compilation time is excluded. Probability
calculations use float64 arithmetic and an 80\% observation mask. Loss timings include backward.

{\small\def\LTcaptype{none} % do not increment counter
\begin{longtable}[]{@{}lrrr@{}}
\toprule\noalign{}
Operation & CPU fused C++ & CPU eager PyTorch & CPU compiled PyTorch \\
\midrule\noalign{}
\endhead
\bottomrule\noalign{}
\endlastfoot
Multilabel Brier, probabilities & 1.892 ms & 18.446 ms & 7.219 ms \\
Multilabel NLL, probabilities & 6.810 ms & 38.126 ms & 11.034 ms \\
\end{longtable}
}

For squared Dice on this GPU workload, compiling native PyTorch was fastest. Choose the
backend and compilation mode using representative shapes, not the implementation language:

{\small\def\LTcaptype{none} % do not increment counter
\begin{longtable}[]{@{}lrrr@{}}
\toprule\noalign{}
GPU forward/backward & Eager PyTorch & Compiled PyTorch & Eager fused Triton \\
\midrule\noalign{}
\endhead
\bottomrule\noalign{}
\endlastfoot
Squared Dice, sigmoid & 0.657 ms & 0.267 ms & 0.575 ms \\
Lovasz-hinge, independent labels & 1.666 ms & 0.865 ms & Uses PyTorch sorting \\
\end{longtable}
}

Exact per-case voxel AP took 131.745 ms on CPU and 4.293 ms on GPU; 256-bin approximation
took 19.352 ms and 3.620 ms. This compares different ranking definitions, not interchangeable
implementations. Pooling updates can change both values and timing. Binned streaming retains
only histogram state; its current-batch working memory still scales with voxel count.

Three physical volume quantities evaluated together took 1.589 ms CPU / 0.423 ms GPU,
versus 4.644 ms / 0.853 ms for three separate calls. Their shared counts avoid repeated scans.
The implementations are included in the \texttt{v0.0.1a5} snapshot. Reproduce these measurements
from that source checkout:

\begin{lstlisting}[language=bash]
OMP_NUM_THREADS=2 MKL_NUM_THREADS=2 TORCHINDUCTOR_COMPILE_THREADS=1 \
  python benchmarks/bench_probability.py --device cpu --size 64 --compiled --output local/probability-cpu.json
OMP_NUM_THREADS=2 MKL_NUM_THREADS=2 TORCHINDUCTOR_COMPILE_THREADS=1 \
  python benchmarks/bench_completion.py --device cuda --size 64 --compiled --output local/completion-gpu.json
\end{lstlisting}

Both scripts support \texttt{-\/-device\ cpu/cuda}, \texttt{-\/-dim\ 2/3} and configurable sizes. Output records
raw samples, allocation measurements, configuration, environment and source hashes. Different
mask density, class counts, hardware and compiler versions can change the best path.

\section{Instances, thinning, topology and soft-label losses}\label{performance:instances-thinning-topology-and-soft-label-losses}

\texttt{benchmarks/bench\_segmentation\_coverage.py}, in the \texttt{v0.0.1a6} source, times the segmentation additions of 0.0.1a6 on
seeded synthetic data, and first checks that compared paths agree. Medians of 10 steps after 3
warm-up steps, two CPU threads, Ryzen 9 5900XT and RTX 5070, PyTorch 2.14.1, float32 losses:

{\small\def\LTcaptype{none} % do not increment counter
\begin{longtable}[]{@{}
  >{\raggedright\arraybackslash}p{(\linewidth - 4\tabcolsep) * \real{0.5500}}
  >{\raggedleft\arraybackslash}p{(\linewidth - 4\tabcolsep) * \real{0.2000}}
  >{\raggedleft\arraybackslash}p{(\linewidth - 4\tabcolsep) * \real{0.2500}}@{}}
\toprule\noalign{}
\begin{minipage}[b]{\linewidth}\raggedright
Operation
\end{minipage} & \begin{minipage}[b]{\linewidth}\raggedleft
CPU
\end{minipage} & \begin{minipage}[b]{\linewidth}\raggedleft
CUDA
\end{minipage} \\
\midrule\noalign{}
\endhead
\bottomrule\noalign{}
\endlastfoot
\texttt{instances.evaluate}, 1024², 2000 touching nuclei & 16.9 ms & 9.7 ms \\
\texttt{instances.dsb2018}, same input (10 thresholds) & 19.5 ms & 12.1 ms \\
\texttt{skeletonize(method="zhang")}, 1024² vessels (scikit-image: 6.8 ms) & 3.4 ms & 2.7 ms, via the host \\
\texttt{betti\_numbers}, 128³, 2 classes & 25.4 ms & 3.9 ms \\
\texttt{Dice.update}, case or pooled, \texttt{2×96×128²}, 4 classes & 1.8 / 1.8 ms & 0.28 / 0.30 ms \\
\texttt{skeleton\_recall\_target}, \texttt{2×64³} (CPU thinning) & 29.9 ms & --- \\
\texttt{SkeletonRecallLoss} / \texttt{DiceCELoss} forward--backward, \texttt{2×3×64³} & 3.8 / 9.2 ms & 0.32 / 0.55 ms \\
\texttt{DiceLoss} forward--backward, soft targets \texttt{2×4×96³}: semimetric / product (PyTorch) / product (fused) & 40.8 / 37.1 ms / --- & 0.76 / 0.73 / 0.54 ms \\
\end{longtable}
}

Pooled aggregation reuses the per-case count pass and adds no measurable cost. Skeleton Recall
is cheaper than the regional loss it is added to, once targets come from the data loader;
thinning targets on every step would dominate. The soft-label Dice runs on PyTorch only and
costs about 40\% more than the fused product form on this GPU workload. A fused kernel needs a
measured case before it is worth adding. Zhang thinning speed relative to scikit-image depends
on how much foreground there is to thin.

\begin{lstlisting}[language=bash]
OMP_NUM_THREADS=2 MKL_NUM_THREADS=2 TORCHINDUCTOR_COMPILE_THREADS=1 \
  python benchmarks/bench_segmentation_coverage.py --device cuda --output local/segmentation-cuda.json
\end{lstlisting}

\chapter{Other functionality}\label{ch:other}

This manual describes the overlap and surface metrics, the main losses and the conventions
behind them in detail. The package contains further modules, each documented on its own page
of the \href{https://github.com/rodolfo-viana/exacta/wiki}{Wiki}. They follow the same rules:
every function states its reduction and its behaviour for empty inputs, and every convention
is an argument.

\begin{longtable}{@{}
  >{\raggedright\arraybackslash}p{0.30\linewidth}
  >{\raggedright\arraybackslash}p{0.44\linewidth}
  >{\raggedright\arraybackslash}p{0.20\linewidth}@{}}
\toprule
Module or namespace & Content & Wiki page \\
\midrule
\endhead
\texttt{exacta.functional}, \texttt{exacta.metrics} &
Overlap, surface, physical area and volume, probability and calibration, ranking (AP,
AUROC), centreline and Betti-number scores; shared evaluation of many metrics with
\texttt{evaluate} and \texttt{MetricCollection}; stateful classes with distributed
aggregation &
Metrics, Medical-segmentation-2D \\
\texttt{exacta.losses} &
Dice, Tversky, focal, generalized Dice and their compositions; top-$k$ cross-entropy,
Lov\'asz, boundary, Hausdorff distance-transform, soft clDice, Skeleton Recall and
Wasserstein Dice losses; deep supervision and distributed global reductions &
Losses, Advanced-segmentation \\
\texttt{exacta.lesions}, \texttt{exacta.components}, \texttt{exacta.detection} &
Connected components, lesion matching (greedy, Hungarian, merge, overlap) and lesion
detection (AP, FROC) &
Lesion-evaluation \\
\texttt{exacta.instances} &
Evaluation of instance label maps: panoptic, aggregated Jaccard and SEG scores, and the 2018
Data Science Bowl score &
Instance-evaluation \\
\texttt{exacta.functional.}\allowbreak\texttt{classification}, \texttt{exacta.functional.}\allowbreak\texttt{regression} &
Dataset-level classification and regression metrics &
Classification, Regression \\
\texttt{exacta.conventions} &
Reporting signatures that write a score's full definition in one string (Section~\ref{metrics:reporting-signatures-since-0.1.0}) &
Metrics \\
\bottomrule
\end{longtable}

\chapter{The example driver}\label{ch:driver}

\texttt{examples/driver.py} runs a representative set of the public API on small synthetic
inputs and prints the results. \texttt{examples/driver\_expected.txt} holds its output on a CPU,
and \texttt{tests/test\_example\_driver.py} compares the two. The driver is the quickest way to
see that an installation works and to compare it with the expected values.

\section{Running it}\label{driver:running}

From the root of a source checkout, with Exacta installed:
\begin{lstlisting}[language=bash]
python examples/driver.py                              # CPU
python examples/driver.py --device cuda                # the same on a CUDA device
python examples/driver.py 2> /dev/null | diff - examples/driver_expected.txt
\end{lstlisting}
It takes a few seconds, most of it spent importing PyTorch. The standard output is the result;
a line with the versions of Exacta, PyTorch, NumPy and Python, the platform, the device and the
number of threads goes to standard error, so that the standard output does not depend on the
machine. The exit status is zero when every check at the end of the output passes.

\section{What it does}\label{driver:what}

The driver draws no random numbers. Every input is built from integer arithmetic: two cases of
$22\times26\times24$ voxels whose target is a discrete ball and whose predictions are the ball
moved by three voxels along the last axis (plus a small separate ball), and a larger ball. The
spacing is $(1.2, 0.4, 0.4)$ millimetres: anisotropic, and none of its values is exactly
representable in binary, so products such as $3\times0.4$ are inexact in binary64. The sections are:

\begin{enumerate}[nosep]
\item the inputs;
\item the exact distance transform of \texttt{exacta.losses.distance\_maps} on a
  $5\times6\times7$ grid with spacing $(1.3,0.3,0.7)$, compared bit for bit with a brute-force
  search written in the driver;
\item HD, HD95 and ASSD with the three surface models (\texttt{voxels}, \texttt{surfels},
  \texttt{mesh}), and the overlap scores and counts;
\item NSD at a tolerance of 1.2 mm, where distances tie with the tolerance, with the default tie
  margin and with \texttt{tie\_rtol=0};
\item the presets that reproduce MONAI, MedPy, DeepMind's \texttt{surface-distance} and
  MeshMetrics (Chapter~\ref{ch:tools});
\item a stateful metric (\texttt{HD95}) and its \texttt{per\_case}, \texttt{per\_class} and
  \texttt{mean}, and the reporting signature of a score, parsed back and recomputed to the
  same bits;
\item \texttt{DiceLoss} and \texttt{DiceCELoss} with value and gradient norm, and the check that
  the second is the first plus PyTorch's cross-entropy.
\end{enumerate}

Numbers are written with Python's \texttt{repr}, which prints a float64 with the digits needed to
read it back bit for bit. Lines without a marker come from computations that Exacta rounds or
counts exactly; they are the same on every platform, thread count and device. Lines that end
with \texttt{(approx)} are the losses, which run on PyTorch operations or fused Triton kernels, and
ASSD and NSD with mesh weights, which sum triangle areas that PyTorch computes (a cross product
and a norm). The last bits of both depend on the PyTorch build and the device: the mesh ASSD
differed by one unit in the last place on Linux aarch64 with PyTorch 2.15, and not with 2.4.1.
The test compares these lines to a relative tolerance of $10^{-9}$ and everything else as text.

\section{Reading the output}\label{driver:reading}

The output below was produced on x86-64 Linux with PyTorch 2.14.1 on the CPU.
Three lines show what the conventions of Chapter~\ref{ch:conventions} do:
\begin{itemize}[nosep]
\item \texttt{3 * 0.4} is \texttt{1.2000000000000002}, so the 1.2 mm distance of three voxels
  lies above a tolerance of 1.2 in binary64.
\item The voxel NSD of case 0 is \texttt{0.9717} with the default margin and \texttt{0.7778} with
  \texttt{tie\_rtol=0}: the surface elements that lie exactly on the tolerance are counted in
  the first and left out in the second. The surfel and mesh values move in the same way.
\item HD95 of case 0 is \texttt{3.2242} with \texttt{combine="max"} (MONAI), but with
  \texttt{combine="pooled"} (MedPy) it is \texttt{1.2000000000000002}: one name, two quantities.
\end{itemize}

\lstinputlisting[language=output,basicstyle=\ttfamily\scriptsize,caption={Output of
\texttt{python examples/driver.py} on the CPU (\texttt{examples/driver\_expected.txt}).}]{../../examples/driver_expected.txt}

\section{How it was checked}\label{driver:checked}

The output was identical in two runs, and with \texttt{OMP\_NUM\_THREADS} set to 1 and to 8,
including the loss lines. On a CUDA device (\texttt{--device cuda}) every line without the
marker was identical to the file, and the two gradient norms differed from it in the last
digit. The test file checks the same properties: the output equals the expected file, it does
not depend on the thread count, and, when a GPU is present, the exact lines are the same on it.
After a deliberate change of the driver, regenerate the expected file with
\begin{lstlisting}[language=bash]
python examples/driver.py 2> /dev/null > examples/driver_expected.txt
\end{lstlisting}
and review the difference.

\chapter{Tests and the referee subset}\label{ch:tests}

The test suite lives in \texttt{tests/} and runs with \texttt{pytest}. Tests that compare
Exacta with another package skip when the package is missing, and tests that need a GPU skip on
a CPU-only machine; a passing CPU-only run therefore does not establish GPU correctness. The
complete suite takes several minutes on CPU. The \emph{referee subset} selects the tests that
bear on the claims of Chapter~\ref{ch:introduction} and runs in a few minutes on a CPU-only
machine.

\section{An isolated environment}\label{tests:referee}

The commands below build the extension and install the test dependencies into a virtual
environment inside the checkout, so that no other Python environment changes. PyTorch is
installed from its CPU index; use the index for your hardware if you also want the GPU tests.
\begin{lstlisting}[language=bash]
git clone --branch main https://github.com/rodolfo-viana/exacta.git
cd exacta
python -m venv .venv && . .venv/bin/activate
python -m pip install torch --index-url https://download.pytorch.org/whl/cpu
python -m pip install -e ".[test]"
\end{lstlisting}
The editable install compiles the C++ kernels and needs a C++17 compiler (OpenMP is optional).
The reference tools are optional. Without them the tests that use them are skipped, and the
pytest summary counts the skips:
\begin{lstlisting}[language=bash]
python -m pip install monai medpy scikit-image vtk SimpleITK \
  segmentation-models-pytorch \
  "git+https://github.com/google-deepmind/surface-distance.git" \
  "git+https://github.com/gasperpodobnik/MeshMetrics.git"
\end{lstlisting}

\section{Running the referee subset}\label{tests:running}

\begin{lstlisting}[language=bash]
python -m pytest -m referee -q -rs     # the referee subset; -rs lists the skips
python -m pytest -q                    # the complete suite
python examples/driver.py              # the example driver
\end{lstlisting}

\texttt{tests/referee\_subset.txt} lists the tests of the subset as pytest node ids, with
comments that say why each group is there. \texttt{tests/conftest.py} adds the marker
\texttt{referee}, registered in \texttt{pyproject.toml}, to every test that matches an entry:
a file selects all its tests, \texttt{file::test} all parametrizations of a test. No existing
test file is edited. If a listed file is collected but an entry matches nothing, pytest stops
with an error, so a renamed test cannot drop out of the subset unnoticed; the test
\texttt{tests/test\_referee\_subset.py} checks the listing as well.

The subset covers, by group of the listing:
\begin{description}[nosep,leftmargin=0pt,labelindent=0pt]
\sloppy
\item[Exact distance transform.] All of \texttt{test\_edt.py}: agreement with SciPy, equality
  with a brute-force search for the exact floating-point minimum (a property test over random
  masks and spacings), the lines that the exactness argument produces (parabolas that tie far
  apart and within rounding error), extreme spacings, the overflow guard and
  thread-count independence.
\item[Correctly rounded sums.] \texttt{test\_sums.py}: the exact sum equals \texttt{math.fsum}
  in any order, and surface means do not depend on the thread count (on GPU, also on the device).
\item[Metric presets against reference tools.] The surface metrics are compared with NumPy and
  SciPy references and with MONAI, MedPy, DeepMind's \texttt{surface-distance}, MeshMetrics and
  VTK, including NSD ties, in \texttt{test\_surface.py}, \texttt{test\_surfels.py} and
  \texttt{test\_mesh.py}. \texttt{test\_conventions.py} checks the reporting signatures.
  Property tests cover HD, ASSD and NSD, and the overlap and volume scores are compared with MedPy,
  MONAI and SimpleITK.
\item[Losses against compositions.] Independent formulas, MONAI and
  \texttt{segmentation\_models\_pytorch}; fused or composed losses against separately composed
  PyTorch terms, values and gradients, and finite-difference gradient checks.
\item[The example driver.] \texttt{test\_example\_driver.py} (Chapter~\ref{ch:driver}).
\end{description}

\section{Size and runtime}\label{tests:runtime}

On a CPU-only machine with the reference tools of Table~\ref{tab:versions} installed, pytest
collects 2713 tests. The subset selects 614 of them: 581 pass and 33 are skipped because they
need CUDA and Triton, and none is skipped for a missing package. The complete suite runs
2119 tests and skips 594.

Run times, with two OpenMP threads, CUDA hidden, on an AMD Ryzen 9 5900XT (16 cores):
\begin{center}
\small
\begin{tabular}{@{}lr@{}}
\toprule
Run & Test time \\
\midrule
Referee subset, three runs & 27.5, 24.0 and 23.6 s \\
Referee subset, restricted to two cores (\texttt{taskset}) & 25.1 s \\
Referee subset, first measurement under heavy load from other jobs & 70 s \\
Complete suite & 106 s \\
\bottomrule
\end{tabular}
\end{center}
The wall-clock time is about 3 to 5 s longer than the test time reported by pytest (interpreter
start and collection). The numbers depend on the machine, its load and the installed packages;
a slower machine takes proportionally longer. The subset takes about a fifth of the time of the
complete suite because it leaves out the distributed (multi-process), \texttt{torch.compile}
and many input-validation tests, which are slow and do not bear on exactness.

\begin{table}[ht]
\centering
\small
\begin{tabular}{@{}ll@{\qquad}ll@{}}
\toprule
Package & Version & Package & Version \\
\midrule
Python & 3.14.4 & MONAI & 1.6.1 \\
PyTorch & 2.14.1 & MedPy & 0.5.2 \\
NumPy & 2.5.4 & DeepMind \texttt{surface-distance} & 0.1 \\
SciPy & 1.18.1 & MeshMetrics (PyMeshMetrics) & 1.4.1 \\
pytest & 9.1.1 & SimpleITK & 2.5.6 \\
Hypothesis & 6.170.8 & VTK & 9.7.1 \\
scikit-image & 0.26.0 & segmentation-models-pytorch & 0.6.0 \\
\bottomrule
\end{tabular}
\caption{Versions installed when the subset was timed.}
\label{tab:versions}
\end{table}

\section{Checking portability}\label{tests:portability}

\texttt{scripts/portability\_digest.py run} evaluates a fixed set of computations (the
distance transform, connected components, skeletons and the surface metrics with every surface
model) on inputs generated from integer arithmetic, and writes the SHA-256 of each result.
\texttt{compare} checks that the digest files written on different machines agree exactly. This
is the tool behind the statement that the surface scores have the same bits on every platform
and thread count that it was run on.
\begin{lstlisting}[language=bash]
python scripts/portability_digest.py run --out digest-here.json
python scripts/portability_digest.py compare digest-here.json digest-there.json
\end{lstlisting}

\chapter{How to cite}\label{ch:citing}

Cite the software with the version you used, and report the signature of each score
(Section~\ref{metrics:reporting-signatures-since-0.1.0}), which names the version and every
convention behind the number. The following entry describes the software; replace the version
by the one printed by \texttt{exacta.\_\_version\_\_}.
\begin{lstlisting}
@software{viana_exacta,
  author  = {Viana, Rodolfo Orlando},
  title   = {Exacta: fast, reference-validated losses and metrics for deep learning on CPU and GPU},
  version = {X.Y.Z},
  year    = {2026},
  url     = {https://github.com/rodolfo-viana/exacta},
  license = {Apache-2.0},
}
\end{lstlisting}
The source repository is \url{https://github.com/rodolfo-viana/exacta}, the package is published
on PyPI as \texttt{exacta}, and bug reports go to the issue tracker of the repository. Exacta
is licensed under the Apache License 2.0; the licenses of the third-party code that it includes
are in the files \texttt{NOTICE} and \texttt{THIRD\_PARTY\_LICENSES} of the source distribution.
When you compare Exacta with another tool, cite that tool and its version too
(Chapter~\ref{ch:tools}).


\end{document}
